% concatenated from nismot_add-JACK-EDIT-journal-revision.tex
\documentclass[10pt]{amsart}
\usepackage[T1]{fontenc}
\usepackage[margin=1.2in,marginparsep=0.1in,marginparwidth=1in]{geometry}
\usepackage{amssymb,amsmath,amsthm,amstext,amscd,latexsym,graphics,graphicx,bbm,caption}
\usepackage[usenames,dvipsnames,svgnames,table]{xcolor}
\usepackage[plainpages=false,colorlinks=true, pagebackref]{hyperref}
\usepackage{tikz}
\usetikzlibrary{cd,positioning}



% set arrows as stealth fighter jets
\tikzset{>=stealth}
\hypersetup{citecolor=Sepia,linkcolor=blue, urlcolor=blue}

\usepackage{cite} \def\citeleft{\begingroup\rm[} \def\citeright{]\endgroup} %upright parenthesis for cite


\makeatletter
\def\@tocline#1#2#3#4#5#6#7{\relax
  \ifnum #1>\c@tocdepth % then omit
  \else
    \par \addpenalty\@secpenalty\addvspace{#2}%
    \begingroup \hyphenpenalty\@M
    \@ifempty{#4}{%
      \@tempdima\csname r@tocindent\number#1\endcsname\relax
    }{%
      \@tempdima#4\relax
    }%
    \parindent\z@ \leftskip#3\relax \advance\leftskip\@tempdima\relax
    \rightskip\@pnumwidth plus4em \parfillskip-\@pnumwidth
    #5\leavevmode\hskip-\@tempdima
      \ifcase #1
       \or\or \hskip 2em \or \hskip 2em \else \hskip 3em \fi%
      #6\nobreak\relax
    \dotfill\hbox to\@pnumwidth{\@tocpagenum{#7}}\par
    \nobreak
    \endgroup
  \fi}
\makeatother

%%%%%%%%%%%%%%%%%%%%%%%%%%%%%%%%%%%%%%%%%%%%%%%%%%%%%%%%%%%%%%%%%

% \swapnumbers

\newtheorem{intro-thm}{Theorem}[]
\theoremstyle{plain}
\newtheorem{thm}{Theorem}[section]
\newtheorem{theorem}[thm]{Theorem}
\newtheorem{q}[thm]{Question}
\newtheorem{lemma}[thm]{Lemma}
\newtheorem{corollary}[thm]{Corollary}
\newtheorem{proposition}[thm]{Proposition}
\newtheorem{conj}[thm]{Conjecture}
\theoremstyle{definition}
\newtheorem{remark}[thm]{Remark}
\newtheorem{point}[thm]{}
\newtheorem{notation}[thm]{Notation}
\newtheorem{rmks}[thm]{Remarks}
\newtheorem{definition}[thm]{Definition}
\newtheorem{example}[thm]{Example}
\newtheorem{claim}[thm]{Claim}
%\numberwithin{equation}{section}
\newcounter{elno}                % This to number lists



%%%%%%%%%%%%%%%%%%%%%%%%%%%%%%%%%%%%%%%%%%%%%%%%%%%%%%%%%%%%%%%%%%%%%%%%%%%%%%

\newcommand{\rup}[1]{\lceil{#1}\rceil}
\newcommand{\rdown}[1]{\lfloor{#1}\rfloor}
\newcommand{\ilim}{\mathop{\varprojlim}\limits} % inverse limit
\newcommand{\dlim}{\mathop{\varinjlim}\limits}  % direct limit
\newcommand{\surj}{\twoheadrightarrow}
\newcommand{\inj}{\hookrightarrow}
\newcommand{\tensor}{\otimes}
\newcommand{\ext}{\bigwedge}
\newcommand{\Intersection}{\bigcap}
\newcommand{\Union}{\bigcup}
\newcommand{\intersection}{\cap}
\newcommand{\union}{\cup}

%%%%%%%%%%%%%%%%%%%%%%%%%%%%% new new commands :) %%%%%%%%%%%%%%%%
\newcommand{\supp}{{\rm Supp}}
\newcommand{\Exceptional}{{\rm Ex}}
\newcommand{\del}{\partial}
\newcommand{\delbar}{\overline{\partial}}
\newcommand{\boldphi}{\mbox{\boldmath $\phi$}}

%%%%%%%%%%%%%%%%%%%%%%%%%%%%%%%%%%%%%%%%%%%%%%%%%%%%%%%%%%%%%%%%%%%%%%%%%%%%%%

\newcommand{\udiv}{\underline{\Div}}

%%%%%%%%%%%%%%%%%

\newcommand{\Proj}{{\P roj}}
\newcommand{\sEnd}{{\sE nd}}
\newcommand{\mc}{\mathcal}
\newcommand{\mb}{\mathbb}
\newcommand{\an}{{\rm an}} 
\newcommand{\red}{{\rm red}}
\newcommand{\codim}{{\rm codim}}
\newcommand{\Dim}{{\rm dim}}
\newcommand{\rank}{{\rm rank}}
\newcommand{\Ker}{{\rm Ker  }}
\newcommand{\Pic}{{\rm Pic}}
\newcommand{\per}{{\rm per}}
\newcommand{\ind}{{\rm ind}}
\newcommand{\Div}{{\rm Div}}
\newcommand{\Hom}{{\rm Hom}}
\newcommand{\Aut}{{\rm Aut}}
\newcommand{\im}{{\rm im}}
\newcommand{\Spec}{{\rm Spec \,}}
\newcommand{\Sing}{{\rm Sing}}
\newcommand{\sing}{{\rm sing}}
\newcommand{\reg}{{\rm reg}}
\newcommand{\Char}{{\rm char}}
\newcommand{\Tr}{{\rm Tr}}
\newcommand{\Gal}{{\rm Gal}}
\newcommand{\Min}{{\rm Min \ }}
\newcommand{\Max}{{\rm Max \ }}
\newcommand{\Alb}{{\rm Alb}\,}
\newcommand{\Mat}{{\rm Mat}}
%\newcommand{\GL}{{\rm GL}\,}        % For the general linear group
\newcommand{\GL}{{\G\L}}
\newcommand{\Ho}{{\rm Ho}}
\newcommand{\ie}{{\it i.e.\/},\ }
\renewcommand{\iff}{\mbox{ $\Longleftrightarrow$ }}
\renewcommand{\tilde}{\widetilde}
% Skriptbuchstaben
\newcommand{\sA}{{\mathcal A}}
\newcommand{\sB}{{\mathcal B}}
\newcommand{\sC}{{\mathcal C}}
\newcommand{\sD}{{\mathcal D}}
\newcommand{\sE}{{\mathcal E}}
\newcommand{\sF}{{\mathcal F}}
\newcommand{\sG}{{\mathcal G}}
\newcommand{\sH}{{\mathcal H}}
\newcommand{\sI}{{\mathcal I}}
\newcommand{\sJ}{{\mathcal J}}
\newcommand{\sK}{{\mathcal K}}
\newcommand{\sL}{{\mathcal L}}
\newcommand{\sM}{{\mathcal M}}
\newcommand{\sN}{{\mathcal N}}
\newcommand{\sO}{{\mathcal O}}
\newcommand{\sP}{{\mathcal P}}
\newcommand{\sQ}{{\mathcal Q}}
\newcommand{\sR}{{\mathcal R}}
\newcommand{\sS}{{\mathcal S}}
\newcommand{\sT}{{\mathcal T}}
\newcommand{\sU}{{\mathcal U}}
\newcommand{\sV}{{\mathcal V}}
\newcommand{\sW}{{\mathcal W}}
\newcommand{\sX}{{\mathcal X}}
\newcommand{\sY}{{\mathcal Y}}
\newcommand{\sZ}{{\mathcal Z}}
% Sonderbuchstaben mit Doppellinie
\newcommand{\A}{{\mathbb A}}
\newcommand{\B}{{\mathbb B}}
\newcommand{\C}{{\mathbb C}}
\newcommand{\D}{{\mathbb D}}
\newcommand{\E}{{\mathbb E}}
\newcommand{\F}{{\mathbb F}}
\newcommand{\G}{{\mathbb G}}
\newcommand{\HH}{{\mathbb H}}
\newcommand{\I}{{\mathbb I}}
\newcommand{\J}{{\mathbb J}}
\newcommand{\M}{{\mathbb M}}
\newcommand{\N}{{\mathbb N}}
\renewcommand{\O}{{\mathbb O}}
\renewcommand{\P}{{\mathbb P}}
\newcommand{\Q}{{\mathbb Q}}
\newcommand{\R}{{\mathbb R}}
\newcommand{\T}{{\mathbb T}}
\newcommand{\U}{{\mathbb U}}
\newcommand{\V}{{\mathbb V}}
\newcommand{\W}{{\mathbb W}}
\newcommand{\X}{{\mathbb X}}
\newcommand{\Y}{{\mathbb Y}}
\newcommand{\Z}{{\mathbb Z}}
\newcommand{\Sh}{\sS h}
\newcommand{\deltaop}{\Delta^{op}(\sS h(Sm/\mathbf{k}))}
\newcommand{\pdeltaop}{\Delta^{op}(P\sS h(Sm/\mathbf{k}))}
%\newcommand{\psh}{\pi_0^{\text{\tiny pre}}}
\newcommand{\psh}{\pi_0}
%\renewcommand{\k}{\mathbf{k}}

\newcommand{\colim}{{\rm colim \,}}
\newcommand{\hocolim}{{\rm hocolim \,}}
\newcommand{\holim}{{\rm holim \,}}
\newcommand{\DM}[2]{\mathbf{DM}_{#2}^{\mathit{eff}}(#1)}


\let\syn\mathsf
\newcommand{\Step}[1]{\underline{\syn{Step \ {#1}}}}
\newcommand{\scr}{\scriptscriptstyle}
\newcommand{\e}[1]{\mathfrak{elim}{( #1)}}
\newcommand{\elim}[1]{\mathfrak{elim}_{\scr #1}}
\newcommand{\dgu}[1]{{{d}}^{\scr 1}_{\scriptstyle #1}}
\newcommand{\dgv}[1]{{{d}}^{\scr 2}_{\scriptstyle #1}}
\newcommand{\dg}{\syn{deg}}


\newcommand{\fixme}[1]{\textcolor{red}{$\bullet$}\marginpar{\small \textcolor{red}{#1}}}
\newcommand{\edit}[1]{\marginpar{\textcolor{Green}{#1}}}

\newcommand{\hash}{\raisebox{.3ex}{\footnotesize \# }}
\newcommand{\Green}[1]{\textcolor{Green}{#1}}

\usepackage{verbatim}

    \newcommand{\spref}[1]{\href{http://stacks.math.columbia.edu/tag/#1}{#1}}
\input{xy}
\xyoption{all}
%%%%%%%%%%%%%%%%%%%%%%%%%%%%%%%%%%%%%%%%%%%%%%%%%%%%%%%%%%%%%%%%%%%%%%%%%%%%%%%
\begin{document}

\title{On the motivic homotopy type of algebraic stacks}

\author{Neeraj Deshmukh}
\address{Department of Mathematics, Boyd Research and Education Center, University of Georgia, Athens, GA 30602, USA}
\email{neeraj.deshmukh@uga.edu}

\author{Jack Hall}
\address{School of Mathematics \& Statistics\\The University of Melbourne\\Parkville,
  VIC, 3010\\Australia}
\email{jack.hall@unimelb.edu.au}


%\subjclass[2000]{}

%\keywords{}

\date{\today}

\maketitle

\begin{abstract}
  We construct smooth presentations of algebraic stacks that are local epimorphisms in the Morel-Voevodsky $\A^1$-homotopy category. As a consequence, we show that the motive of a smooth stack (in Voevodsky's triangulated category of motives) has many of the same properties as the motive of a smooth scheme.
\end{abstract}
\section{Introduction}

An important technique in motivic homotopy theory of algebraic stacks is reduction to the scheme case by means of homotopical descent. This is possible, for instance, when the stacks in question are Nisnevich locally quotient stacks. The results in \cite{AHR19} (and further generalisation in \cite{ahhlr21}) show that stacks with linearly reductive stabilisers are Nisnevich locally quotient stacks. 

In this note, we establish a certain homotopy descent result for any
algebraic stack.  This allows us to conclude that the various
formalisms of motives and motivic homotopy theory\footnote{In the case
  of stacks one has to distinguish between \textit{genuine} vs
  \textit{Kan-extended} motivic spectra. For instance $K$-theory is
  \textit{genuine} but Chow groups are \textit{Kan-extended}. In this
  note we will be concerned with the latter kind of objects.} produce
the correct results for algebraic stacks. To wit, we will show that
the motive of a smooth algebraic stack has similar properties as the
motive of a smooth scheme. This generalises the results in
\cite{cdh20} that were established for Nisnevich locally quotient
stacks. Another improvement, albeit minor, that we can make is to show
that for algebraic stacks with separated diagonal the existing notions
of stable homotopy category coincide: in \cite{chow21}, it is shown
that the stable motivic homotopy category of \cite{chow21} and the
lisse-extended category of \cite{khan2021generalized} are equivalent
when the stack admits a smooth presentation with a Nisnevich local
section (i.e., the \emph{smooth-Nisnevich coverings} of
\cite{pirisi2018}). In Appendix \ref{App:smooth-nis-covers}, we will
show that all algebraic stacks admit such a presentation (Theorem
\ref{MT}).

%In this short note we will discuss some aspects of motivic homotopy theory of algebriac stacks. In \cite{cdh20, kr21, chow21} this topic is discussed albeit with many conditions on the algebraic stacks. In \cite{cdh20} only stacks with Nisnevich local quotient structure in the style of Alper-Hall-Rydh \cite{AHR19, ahhlr21} are considered; while in \cite{chow21} the stable movitic homotopy category is defined for algebraic stacks admitting certain special type of coverings that admit local sections; the lisse-extended theory of \cite{kr21} formalises this construction to any algebraic stack. The stable motivic homotopy category of \cite{chow21} and the lisse-extended category of \cite{kr21} are known to be equivalent when the stack admits a smooth presentation with a Nisnevich local section. 



\begin{definition}
	A \emph{smooth-Nisnevich covering} of algebraic stacks is a morphism $f:\sY\rightarrow \sX$ such that $f$ is smooth and every morphism $\Spec K \rightarrow \sX$ from the spectrum of a field $K$ lifts to $\sY$. % We say that $f$ is  \emph{smooth-Nisnevich covering} when $\sY$ is an algebraic space.
\end{definition}

When $f$ is \'{e}tale and $\sY$ and $\sX$ are algebraic spaces, these
are just the standard Nisnevich coverings.

The existence of smooth-Nisnevich coverings for algebraic stacks goes back to
\cite[\S6]{LMB} (quasi-compact and separated diagonal) and
\cite[Appendix B]{cesnavicius15} (diagonal has relatively separated
fibers). If the stack is of finite type over an infinite field with
affine stabilizers, then it was shown in \cite{pirisi2018} that there
exists a smooth-Nisnevich covering of finite type. In the following theorem, we generalize all of these results by
eliminating separation hypotheses for the existence result and
noetherian hypotheses in the boundedness result; we also prove an
analogous result for Deligne--Mumford stacks. 

\begin{theorem}\label{MT}
	If $\mathcal{X}$ is an algebraic stack, then there exists a
	smooth-Nisnevich covering $p\colon X \to \mathcal{X}$, where $X$ is
	a scheme. If $\mathcal{X}$ is quasi-compact and quasi-separated with
	affine stabilizers, then we may take $X$ to be affine. Moreover, if
	$\mathcal{X}$ is Deligne--Mumford, then we may take $p$ to be
	\'etale.
\end{theorem}



It
was brought to our attention during the preparation of our results
that the existence result was recently proved in a similar way in
\cite{chowdhury2024nonrepresentablesixfunctorformalisms}. However, the qcqs and affine stabilzer part of the theorem is new, and we expect it to have applications beyond motivic homotopy theory (see Remark \ref{remark-mathur-kresch} for one such application).

% Smooth-Nisnevich coverings were first discussed in \cite{pirisi2018}, where it is shown that such coverings exist for finite type stacks with affine stabilisers that are defined over an infinite field. The aim of this note is extend this result to all finite type stacks with affine stabilisers. %In fact, we will show that with some modification the approach of \cite{pirisi2018} extends to any Noetherian base scheme. 
% The following is our main geometric result.

% \begin{theorem}\label{theorem-main}
% 	Let $\sX$ be a quasi-separated algebraic stack with separated diagonal over a scheme $S$. 
% 	\begin{enumerate}
% 		\item There exists a smooth-Nisnevich covering $X\rightarrow \sX$ with $X$ a scheme. 
% 		\item Assume $S$ is Noetherian and that $\sX$ is of finite type over $S$ with affine stabilisers, then $X$ can also be chosen to be of finite type over $S$.
% 	\end{enumerate}
% \end{theorem}

% Note that the argument for Part (1) of the above theorem is essentially contained in \cite[\S 6.7]{LMB}. Part (2) is new in the mixed charactersitic setting and is built from a modification of Pirisi's argument \cite[Proposition 3.6]{pirisi2018} which works over an infinite field.
% %The argument presented here are essentially contained in \cite{LMB} (following \cite{knutson}) and \cite{pirisi2018}. The only difference is that we slightly modify the proof of \cite{pirisi2018} to get the proof for over any Noetherian base.  %The above theorem has the following corollary which makes several homotopical descent arguments accessible for algebraic stacks.

The existence of such covers has implications for motivic homotopy theory of algebraic stacks. More specifically, it shows that the Nisnevich homotopy type of an algebraic stack can be described by a simplicial scheme (or simplicial algebraic space). This makes many homotopical descent arguments accessible for algebraic stacks.


\begin{theorem}\label{corollary-smooth-nisnevich-presentation}
	Let $p:X\rightarrow \sX$ be a smooth-Nisnevich covering over a field $k$. Let $X_{\bullet}$ denote the \v{C}ech nerve of $p$. Then the morphism $p_{\bullet}:X_{\bullet}\rightarrow \sX$ induces an equivalence in the Morel-Voevodsky $\A^1$-homotopy category, $\sH(k)$.
\end{theorem}

A consequence of the above result is that the motive of a smooth algebraic stack continues to enjoy the same properties as the motive of a smooth scheme (see Section \ref{section-applications}). So far this was only known for stacks which satisfied local structure theorems in the sense of \cite{AHR19, ahhlr21} or in a different direction for motives of smooth stacks in the \'{e}tale topology. We describe various results of this kind in the text.

\begin{remark}\label{remark-mathur-kresch}
Another application of Theorem \ref{MT} is the following: In \cite{km22}, the authors use the boundedness result in Theorem \ref{MT} to apply Elkik's approximation technique to the Picard stack and study the following question of Grothendieck: when is the map
\[H^2(X,\G_m)\rightarrow \ilim H^2(X_n,\G_m)\]
injective for a proper morphism $X\rightarrow\Spec A$, where $(A,\mathfrak{m})$ is a complete Noetherian ring and $X_n:= X\times_A \Spec A/\mathfrak{m}^{n+1}$ is the $n$-th infinitesimal thickening.
\end{remark}


%We will present a proof of Theorem \ref{theorem-main} in Section \ref{section-main}, and then describe some applications to motivic cohomology in the final section.\\

\noindent\textbf{Conventions.} We work with algebraic stacks in the sense of \cite{stacks-project}; that is, without separation assumptions unless noted otherwise.\\

\noindent\textbf{Acknowledgements.} Neeraj Deshmukh would like to thank Roberto Pirisi and Tuomas Tajakka for helpful discussions. He also thanks Amit Hogadi,  Siddharth Mathur, Kestutis Cesnavicius and Andrew Kresch for their comments on this note and Utsav Choudhury for teaching him about motives with compact support.\\
The authors would like to thank the referee for their comments.\\
N.D. was supported by the Swiss National Science Foundation (SNF), project 200020\_178729 during the course of this work. He also acknowledges the support of the University of Zurich under the UZH Postdoc Grant, Verf\"{u}gung Nr. FK-22-111.\\
Jack Hall was partially supported
by the Australian Research Council DP210103397 and FT210100405.

%In this short note we will describe the motivic cohomology for any quasi-separated stack. This note is a continuation of the results in [CDH20], and extends the results in that article to any quasi-separated stack.


% \section{Smooth-Nisnevich coverings}
% In order to prove Theorem \ref{theorem-main}, we will use the following results from \cite[\S 6]{LMB}


% \subsection{Moduli of finite \'{e}tale covers}

% Fix a base scheme $S$. For a separated and representable $S$-morphism $\pi:\sX \rightarrow \sY$ of algebraic stacks, we recall the construction and basic properties of the stacks $SEC_d(\sX/\sY)$ and $ET_d(\sX/\sY)$.

% Consider the $d$-fold fibre product $(\sX/\sY)^d:=\sX \times_\sY \sX \times_\sY \ldots \times_\sY \sX$. For any $S$-scheme $U$, an object $u: U\rightarrow (\sX/\sY)^d$ is equivalent to the data of a morphism $U\rightarrow \sY$ together with $d$ sections $x_1,\ldots, x_d$ of the morphism $\sX\times_\sY U \rightarrow U$. Since $\pi$ is separated, the locus where any two such section agree is closed. Thus, the complement of this locus is an open substack which is denoted by $SEC_d(\sX/\sY)$. 

% The symmetric group $S_d$ acts naturally on $(\sX/\sY)^d$ and the restriction of this action to $SEC_d(\sX/\sY)$ is free. We denote the quotient stack by

% \[ET_d(\sX/\sY):= [SEC_d(\sX/\sY)/S_d].\]

% Almost by construction $SEC_d$ and $ET_d$ satisfy the following properties:

% \begin{proposition}[{\cite[Propositions 6.6.2, 6.6.3]{LMB}}]
% 	Let $\pi: \sX\rightarrow\sY$ be a representable and separated morphism of algebraic stacks over $S$. Then
% 	\begin{enumerate}
% 		\item The constructions $SEC_d(\sX/\sY)$ and $ET_d(\sX/\sY)$ behave well with respect to base change. That is, given any morphism $U\rightarrow \sY$, we have cononical isomorphisms $SEC_d(\sX/\sY)\times_\sY U \simeq SEC_d(\sX\times_\sY U/ U)$ and $ET_d(\sX/\sY)\times_\sY U \simeq ET_d(\sX\times_\sY U/ U)$.
% 		\item The natural morphisms $SEC_d(\sX/\sY)\rightarrow \sY$ and $ET_d(\sX/\sY)\rightarrow \sY$ are representable and separated. Moreover, they are smooth (respectively, \'{e}tale) if $\pi$ is.
% 		\item If $\sX$ is an algebraic space, then $SEC_d(\sX/\sY)$ is an algebraic space and, therefore, $ET_d(\sX/\sY)$ is a Deligne-Mumford stack.
% 		\item For any $S$-scheme $U$, a $U$-point of $ET_d(\sX/\sY)$ is equivalent to a morphism $U\rightarrow \sY$ together with a closed subscheme $Z\subset U\times_\sY \sX$ such that $Z\rightarrow U$ is finite \'{e}tale of degree $d$.
% 		\item If $\pi$ is finite \'{e}tale of degree $d$, then $ET_d(\sX/\sY)\simeq \sY$.
% 	\end{enumerate}
% \end{proposition}


% \subsection{Proof of Theorem \ref{theorem-main}}\label{section-main}

% For part (2) of the theorem, we will need the following result about stacks admitting stratification by global quotient stacks as found in \cite{KreschChow}.

% \begin{definition}
% 	Let $\sX$ be an algebraic stack over $S$. We say that $\sX$ admits a stratification by global quotient stacks if there exists a finite chain $\sX\supset \sX_1\supset \ldots \sX_n$ of closed substacks such that for each $i$, $\sX_i \setminus \sX_{i-1}$ is a global quotient stack over $S$.
% \end{definition}


% \begin{theorem}\cite[Theorem 3.5.9]{KreschChow}\label{theorem-kresch-stratification}
% 	Let $\sX$ be an algebraic stack of finite type over a Noetherian scheme $S$. Assume that $\sX$ has affine stabilizers at geometric points. Then $\sX$ admits a stratification global quotient stacks over $S$.
% \end{theorem}


% We also need the following result whose proof can also be found in \cite[Th\'{e}or\`{e}me 6.1]{LMB}

% \begin{proposition}
% 	Let $\sY$ be a Deligne-Mumford stack. Assume that $\sY$ is of the form $[Y/G]$ with $Y$ an algebraic space and $G$ a finite group. Then there exists a smooth covering $X\rightarrow \sY$ such that given any semi-local ring $\sO$ any $\sO$-valued point of $\sY$ lifts to $X$.
% \end{proposition}
% \begin{comment}
% \begin{proof}
% 	Choose a faithful representation $G \hookrightarrow GL_n$, and take $X:= Y\times^{G} GL_n$. Then $\sY\simeq [X/GL_n]$ and the natural map $ X \rightarrow \sY$ has the desired property.
% \end{proof}
% \end{comment}

% We will now prove Theorem \ref{theorem-main}.

% %To prove this theorem we first show that there exists a smooth-Nisnevich representable morphism $\sY\rightarrow \sX$ with $\sY$ a Deligne-Mumford stack. Moreover, we can arrange that $\sY$ is the quotient of an algebraic space by a finite group. This reduction reduces the proof to the following


% \begin{proof}[Proof of \ref{theorem-main}]
% 	For both part (a) and part (b), we will show that we can find such a covering with $X$ an algebraic space. As any algebraic space admits a Nisnevich covering by a scheme (see \cite[Theorem 6.4]{knutson}),we conclude by replacing $X$ with a Nisnevich covering by a scheme.
	
% 	\noindent\underline{Proof of part (1).}
% 	The argument for (1) is essentially in \cite[\S 6.7]{LMB} and we reproduce it here for completeness.
% 	It suffices to consider the case when $\sX$ is quasi-compact.
% 	Let $Y\rightarrow \sX$ be a smooth presentation with $Y$ affine.  Let $K$ be a field and choose a $K$-valued point $\Spec K\rightarrow \sX$. Then there exists a finite separable extension $L\supset K$ and a closed immersion $f: \Spec L \rightarrow Y$. Let $d:= [L:K]$. Denote by $ET_d(Y/\sX)$ the moduli of degree $d$ \'{e}tale covers. Then the closed immersion $f: \Spec L \rightarrow X$ defines a $K$-point of $ET_d(XY/\sX)$ giving us a commutative diagram,
% 	\begin{center}
% 	\begin{tikzcd}[cramped,column sep= small]
% 	 & ET_d(Y/\sX)\arrow[d]\\
% 	\Spec K \arrow[ur]\arrow[r] & \sX	
% 	\end{tikzcd}
% 	\end{center}

% 	As $ET_d(Y/\sX):= [SEC_d(Y/\sX)/S_d]$ and $Y$ is affine, $SEC_d(Y/\sX)$ is an algebraic space. Applying the previous proposition, we see that $\Spec K$ lifts to $X_d:= SEC_d(Y/\sX)\times^{S_d} GL_n$. Then the countable collection $\{X_n\}_n$ gives a smooth-Nisnevich covering $\cup_n X_n\rightarrow \sX$. As every algebraic space admits a Nisnevich covering by a scheme we can also assume that the family $\{X_n\}$ is a countable collection of schemes.\\
	
% 	\noindent\underline{Proof of part (2).} Assume $\sX$ is finite type with affine stabilisers and that $S$ is Noetherian. Let $\sA:=\{X_n\}_n$ be the countable family constructed in part (1), and $p:\cup_n X_n \rightarrow \sX$ be the smooth-Nisnevich covering. By the Theorem \ref{theorem-kresch-stratification}, $\sX$ admits a stratification by global quotient stacks. Thus, it suffices to show that for each stratum $\sY_i:=\sX_i\setminus\sX_{i-1}$, there is a finite collection $X_{i_1}, X_{i_2},\ldots, X_{i_k} \in \sA$ such that $\cup_{k}(X_{i_k}\times_\sX \sY_i) \rightarrow \sY_i$ is a smooth-Nisnevich covering. The union of all such $X_{i_k}$ is a finite type scheme that is a smooth-Nisnevich covering of $\sX$.
% 	%Further, since any morphism $\Spec K\rightarrow \sX$ factors through the residuel gerbe, if $\Spec K\times_\sX [X/GL_n] = \emptyset$ then it factors through the closed substack $\sZ$. 
% 	Thus, it suffices to consider the case of global quotient stacks. 
	
% 	Let $[Y/GL_n]$ be a global quotient stack. Let $\cup_n X_n \rightarrow [Y/GL_n]$ be the countable family of part (1). We have the following cartesian diagram where each arrow is a smooth-Nisnevich covering:
% 	\begin{center}
% 	\begin{tikzcd}
% 		\cup_n (X_n \times_{[Y/GL_n]} Y) \arrow[r]\arrow[d] & Y\arrow[d] \\
% 		\cup_n X_n\arrow[r] & {[Y/GL_n]}
% 	\end{tikzcd}
% 	\end{center}
% 	Denote $\cup_n (X_n \times_{[Y/GL_n]} Y):= \cup_n Z_n$. Since $\cup_n Z_n\rightarrow Y$ is smooth-Nisnevich, the generic point of $Y$ admits a lift to some $Z_k$. Thus, there is an open subset $U\subset Y$ that lifts to $Z_k$. We may now consider the complement of $U$, and apply the same procedure. This gives a descending sequence $Y\subset Y_1 \subset \ldots$ whose generic points lift to $\cup_n Z_n$. Since $Y$ is Noetherian, there exists an $N$ such that $\cup_{n\leq N}Z_n$ stabilises. 
	
% 	We will show that the map $ \cup_{n\leq N}Z_n\rightarrow Y$ is a smooth-Nisnevich covering. Let $x:\Spec k\rightarrow Y$ be a morphism. By construction, there exists an $i$ such that the morphism $x$ factors as $\Spec k\rightarrow Y_i\setminus Y_{i-1}\subset Y$. Then the lifting $Y_i\setminus Y_{i-1}\subset Z_i\subset \cup_{n\leq N}Z_n$, gives a lift of $x: \Spec k\rightarrow Y$. Further, as $Y\rightarrow [Y/GL_n]$ is a smooth-Nisnevich covering, it is easy to see that $\cup_{n\leq N} X_n \rightarrow [Y/GL_n]$ is also a smooth-Nisnevich covering.
% \end{proof}


% \begin{remark}
% 	We note that the proof of part (2) above is essentially the same argument as in \cite[Proposition 3.6]{pirisi2018}. A careful choice of presentations for the strata allows us to %The only distinction is that we use the covering $X\rightarrow [X/GL_n]$ instead of a schematic open in a vector bundle as in \cite{pirisi2018}. This allows us 
% 	to relax the characteristic zero assumption in \cite[Proposition 3.6]{pirisi2018} and also deal with the mixed characteristic situation.
% \end{remark}

% \begin{remark}
% 	The argument for part (1) is extracted from the proof of \cite[Theorem 6.7]{LMB} which states for any field $k$ and a $k$-point of $\sX$ there exists smooth affine neighbourhood that lifts the given $k$-point. While it is tempting to take union over all such neighbourhood to get a smooth-Nisnevich covering, it is unclear to me how to do this directly. This is due the fact that there is no notion of a residue field for algebraic stacks, hence there is no concept of ``minimal" fields through which any field valued point factors.
% \end{remark}
% In fact smooth-Nisnevich coverings also lift Hensel local points by the following corollary.

% \begin{corollary}\label{corollary-lift-hensel-local-points}
% 	Let $X\rightarrow\sX$ be a smooth-Nisnevich covering. Let $\Spec \sO$ be the spectrum of a Henselian local ring $\sO$. Then any morphism $\Spec \sO\rightarrow \sX$ lifts to a morphism $\Spec \sO\rightarrow X$.
% \end{corollary}
% \begin{proof}
% 	Let $k$ be the residue field of $\sO$ and $X_\sO:= X\times_\sX \Spec \sO$ denote the base change of $X$ to $\Spec \sO$. Then the composite $\Spec k \rightarrow \Spec \sO \rightarrow \sX$ lifts to a morphism $\Spec k \rightarrow X$. This also gives us a morphism $\Spec k \rightarrow X_\sO$. Now, as $X_\sO \rightarrow \Spec \sO$ is smooth, we have a surjection of sets $X_\sO(\sO)\twoheadrightarrow X_\sO(k)$. Thus, we have a section $\Spec \sO \rightarrow X_\sO$. Composing with the projection map $X_\sO \rightarrow X$ gives us the desire lift.
% \end{proof}

% \subsection{The non-separated case}
% In general, the proof of theorem \ref{theorem-main}(1) breaks down in the non-separated situation. However, we will describe a particular situation where the statement hold. This is essentially based on \cite[Appendix B]{cesnavicius15}. We make the following definition

% \begin{definition}
% 	Let $f:\sX\rightarrow \sY$ be a morphism of algebraic stacks defined over a scheme $S$. We say that $f:\sX\rightarrow\sY$ is $S$-fibrewise separated, if for every $s\in S$, the fibre $\sX_s\rightarrow \sY_s$ is separated.
% \end{definition}


% \begin{proposition}
% 	Let $\sX$ be an algebraic stack over a scheme $S$. Suppose that the diagonal $\Delta_{\sX/S}$ is $S$-fibrewise separated. Then there exists a smooth-Nisnevich morphism $f:X\rightarrow\sX$ with $X$ a scheme.
% \end{proposition}

% \begin{proof}
% 	The proof essentially follows from \cite[Appendix B]{cesnavicius15}. One just needs to do careful book-keeping to ensure that the end result is a set.
	
% 	We will briefly sketch the steps of the argument (see \cite[Appendix B]{cesnavicius15} for details):
	
% 	Step 1: Let $X\rightarrow \sX$ be a smooth atlas. Consider the $n$-fold product $(X/\sX)^n$ with the action of $S_n$ given by permuting the coordinates. The associated quotient stack  $[(X/\sX)^n/S_n]\rightarrow \sX$ is smooth over $\sX$ with separated diagonal. The mophism $[(X/\sX)^n/S_n]\rightarrow \sX$ behaves well with respect to base change. Let $Y_n\rightarrow [(X/\sX)^n/S_n]$ be a smooth-nisnevich covering. We will show that $\sqcup Y_n\rightarrow \sX$ has the required property. Moreover, note that it suffices to prove that any $k$-point of $\sX$ lifts to $[(X/\sX)^n/S_n]$ for some $n$.
	
% 	Step 2: Let $x:\Spec k\rightarrow\sX$ be a $k$-point. Let $s:\Spec k \rightarrow S$ be its image in $S$, and consider the fibre of $\sX$ over $s$, $\sX_s\rightarrow\Spec k$. As $\Delta_{\sX/S}$ is fibrewise separated, the morphism $\Spec k\rightarrow \sX_s$ admits a lift to $Et_n(X_s/\sX_s)$ for some $n$. Further, $Et_n(X_s/\sX_s)\rightarrow [(X/\sX)^n/S_n]_s$ is an  open immersion, giving us a map, $\Spec k \rightarrow [(X/\sX)^n/S_n]_s\rightarrow[(X/\sX)^n/S_n]$, as required.
% \end{proof}

% \begin{comment}
% The separated diagonal case:

% Let $\sX$ be an algebraic stack, and $x:\Spec k\rightarrow \sX$ be a $k$-valued point. As $\sX$  can be writted as a union of quasi-compact open substacks, we may assume that $\sX$ is quasi-compact. Let $X\rightarrow \sX$ be a smooth presentation with $X$ affine. Let $X_k$ be the fibre over the $k$-point $x$. $X_k$ is a smooth separated algebraic space over $k$. Thus, there exists a closed immersion $\Spec L\rightarrow X_k$ such that $L/k$ is a finite separable extension. Thus, the $k$-point $x$ lifts to a point $\Spec k\rightarrow Et_n(X/\sX)$ for some $n$. We have an identification $Et_n(X/\sX)\simeq [Sec_n(X/\sX)/S_n]$, where $Sec_n(X/\sX)$ is an algebraic space not necessarily quasi-separated. Write $Y:= Sec_n(X/\sX) $ Thus, we get a lift $y:\Spec k\rightarrow Y$ of $x$.

% To get a smooth-Nisnevich covering of $X$, we need to construct a covering of $Y$ which lifts the $k$-point $y$. As $Y$ is may not be quasi-separated, it may not admit a Nisnevich covering unlike \cite[6.7]{knutson}. In fact, the map $\Spec k\rightarrow Y$ may not even be a monomorphism. We resolve this situation by essentially repeating the above argument.

% As before, we can replace $Y$ with a quasi-compact open. Let $V\rightarrow Y$ be affine \'{e}tale presentation. Then we have a lift of $y$ to a point of $Et_m(V/Y)$ for some $m$. Further, B.3(a) tells us that $Et_m(V/Y)$ is a separated algebraic space. Thus, we have a Nisnvich covering of $Et_m(V/Y)$. Taking the union over a covering family of $\sX$ and then $Y$ gives us smooth-Nisnevich covering $W\rightarrow\sX$ as required.

% \end{comment}

% \begin{remark}
% 	The covering constructed above can be very large. We have described this construction more as an idle curiosity than anything else. In practice, the case of qcqs algebraic stacks with separated diagonal suffices for most applications.
% \end{remark}


\section{Smooth-Nisnevich covers}\label{App:smooth-nis-covers}
% Recall that a morphism of algebraic stacks $p \colon X \to \mathcal{X}$ is a
% \emph{smooth-Nisnevich covering} if it is smooth and for every
% morphism $x \colon \Spec k \to \mathcal{X}$, where $k$ is a field,
% there is a lift $\tilde{x} \colon \Spec k \to X$ of
% $x$. 
% Smooth-Nisnevich coverings were used by Pirisi \cite{MR3739195}
% and Deshmukh \cite{deshmukh2023motivichomotopytypealgebraic} to define
% cohomological and motivic invariants for algebraic stacks; also see
% \cite{MR3482265} and
% \cite[\S6]{MR1771927}. 
%The main result of this section is the following.
% \begin{remark}
%   Note that Theorem \ref{MT} implies that if $\mathcal{X}$ is an
%   algebraic stack with affine-stabilizers of finite type over a field $k$, then
%   there is a smooth-Nisnevich cover $\Spec A \to \mathcal{X}$, where
%   $A$ is a finitely generated $k$-algebra.
% \end{remark}
% If the inertia of $\mathcal{X}$ has relatively separated fibers, then the existence part of Theorem \ref{MT} was already known \cite[Thm.~B.5]{cesnavicius15}.
% Other ????  separated inertia separated diagonal over a field in
% \cite[Prop.~3.6]{MR3739195} and if $\mathcal{X}$ has separated
% diagonal in
% \cite[Thm~1.2]{deshmukh2023motivichomotopytypealgebraic}. 
% The boundedness in Theorem \ref{MT} was recently shown for stacks of finite type over a
% noetherian base with affine stabilizers in
% \cite{deshmukh2023motivichomotopytypealgebraic}. In particular, we
% remove all separation assumptions and noetherian hypotheses. This is
% accomplished using techniques from \cite{MR2821738,etale_dev_add}, the
% latter of which were developed to solve similar problems in the
% Nisnevich topology. The boundedness argument we use is similar in
% spirit to that in
% \cite{deshmukh2023motivichomotopytypealgebraic,MR3739195}, but is
% simplified with our Proposition \ref{P:boundedness} and splitting
% covers.
% \subsection*{Acknowledgements}
% The author would like to thank Neeraj Deshmukh and Lance Gurney for
% helpful comments and suggestions.
%\section{Smooth-Nisnevich coverings}

Nisnevich coverings for algebraic stacks were discussed in
\cite[\S3]{MR3754421}. The difference between smooth-Nisnevich and
Nisnevich coverings is that the latter morphisms are \'etale and
require a lift of the residual gerbe---not just a rational point. Like
their Nisnevich counterparts, however, smooth-Nisnevich coverings are
clearly stable under composition and base change. We begin with some
examples.
\begin{example}\label{ex:alg-sp}
	Let $f \colon X \to S$ be a smooth morphism of stacks, where $S$ is
	a quasi-separated algebraic space. Then every point $s\in |S|$ has a
	residue field $\kappa(s)$ \cite[Tag
	\spref{03JV}]{stacks-project}. Hence, $f$ is a smooth-Nisnevich
	covering if and only if for every $s \in |S|$, the morphism
	$f_s \colon X\times_S \Spec \kappa(s) \to \Spec \kappa(s)$ admits a
	section. In particular, if $X$ is also an algebraic space, then
	\'etale smooth-Nisnevich coverings are equivalent to Nisnevich
	covers in the sense of \cite[\S3]{MR3754421}. Thus, if
	$\mathcal{X}$ is a quasi-separated algebraic space, then Theorem
	\ref{MT} holds for $\mathcal{X}$ \cite[Prop.~5.7.6]{MR308104}.
\end{example}
\begin{example}\label{ex:gln-nis-sm}
	$\Spec \Z \to B\mathrm{GL}_{n,\Z}$ is a smooth-Nisnevich covering. More
	generally, if $Y \to S$ is a morphism of quasi-separated algebraic
	spaces such that $Y$ admits an $S$-action of $\mathrm{GL}_{n,S}$, then
	Theorem \ref{MT} holds for the quotient stack
	$[Y/\mathrm{GL}_{n,S}]$. Indeed, passing to a Nisnevich covering of $S$, we
	may assume that $S$ is an affine scheme. Then $Y \to [Y/\mathrm{GL}_{n,S}]$
	is a smooth-Nisnevich morphism (it is the base change of
	$\Spec \Z \to B\mathrm{GL}_n$). Passing to a Nisnevich \'etale cover of $Y$
	finishes the argument. This holds more generally for $[Y/G]$, where
	$G \to S$ is a smooth group algebraic space such that
	$\mathrm{H}^1((\Spec K)_{\mathrm{et}},G_{\Spec K})=0$ for all
	$\Spec K \to S$.
\end{example}
\begin{example}\label{ex:quasi-dm}
	If $\mathcal{X}$ is an algebraic stack with quasi-finite diagonal,
	then Theorem \ref{MT} holds for $\mathcal{X}$. We may of course
	assume that $\mathcal{X}$ is quasi-compact. By
	\cite[Thm.~4.1]{MR3754421} there exist morphisms of algebraic
	stacks $V \xrightarrow{v} \mathcal{W} \xrightarrow{w} \mathcal{X}$
	such that $V$ is an affine scheme, $v$ is finite and faithfully flat
	of finite presentation, and $w$ is a Nisnevich \'etale
	covering. Hence, we may replace $\mathcal{X}$ by $\mathcal{W}$ and
	assume that $\mathcal{X}$ admits a finite and faithfully flat cover
	of finite presentation by an affine scheme $V$. By
	\cite[Prop.~4.3(vii)]{GrossRes}, $\mathcal{X}$ has the
	resolution property and so the Totaro--Gross Theorem
	\cite[Thm.~5.4]{GrossRes} implies that
	$\mathcal{X} \simeq [X/\mathrm{GL}_n]$, where $X$ is quasi-affine. Now apply
	Example \ref{ex:gln-nis-sm}.
\end{example}
Let $f \colon X \to S$ be a morphism of algebraic stacks. A
\emph{splitting cover} for $f$ is a sequence of quasi-compact immersions
$Z_j \hookrightarrow S$ for $j=1$, $\dots$, $r$ with
$|S|=\cup_{j=1}^r |Z_j|$ such that
$f_{Z_j} \colon X\times_S Z_j \to Z_j$ admits a section for each $j$.
\begin{remark}\label{R:splitting-sequence}
	If $S$ is quasi-compact, then the existence of a splitting cover
	for $f$ is equivalent to the existence of a \emph{splitting
		sequence} for $f$; that is, there is a filtration
	$\emptyset = S_0 \subseteq S_1 \subseteq \cdots \subseteq S_t = S$
	by quasi-compact open subsets such that
	$f_{S_{i}\setminus S_{i-1}} \colon X \times_S (S_{i}\setminus
	S_{i-1}) \to S_i\setminus S_{i-1}$ admits a section for each $i=1$,
	$\dots$, $t$. The equivalence is due to the close relationship
	between partitions and filtrations in the constructible topology
	\cite[Tag \spref{09XY}]{stacks-project}.
\end{remark}
% sequence of quasi-compact open immersions
% \[
%   \emptyset = S_0 \subseteq S_1 \subseteq \cdots \subseteq S_r = S
% \]
% such that $f$ restricted to each strata $Z_i = S_i\setminus S_{i-1}$, when given the
% reduced structure,
%admits a section for each $i=1$, $\dots$, $r$.
% \begin{remark}\label{R:strata-filtrations}
%   Let $S$ be a quasi-compact and quasi-separated algebraic stack.  If
%   there is a partition $|S| = \amalg_{j=1}^t |Z_j|$, where the
%   $Z_j \hookrightarrow S$ are finitely presented immersions of
%   algebraic stacks, then there is a filtration
%   $\emptyset = S_0 \subseteq S_1 \subseteq \cdots \subseteq S_r = S$
%   and each strata $S_{i}\setminus S_{i-1}$ is contained in some $Z_j$
%   for $i=1$, $\dots$, $r$. To see this, we may write each
%   $Z_j = U_j\setminus V_j$, where $V_j \subseteq U_j \subseteq S$ are
%   quasi-compact open subsets. Let $\Sigma = \{U_j,V_j\}_{j=1}^t$ and
%   let $\bar{\Sigma}$ be the smallest collection of subsets of $|S|$
%   containing $\Sigma$ that is closed under finite intersections; then
%   $\bar{\Sigma}$ is finite and every element of
% \end{remark}
The
following lemma is a smooth analog of \cite[Prop.~3.3]{MR3754421}.
\begin{lemma}\label{L:split-covering-sm}
	Let $f \colon X \to S$ be a smooth morphism of algebraic stacks. If
	$S$ is a quasi-compact and quasi-separated algebraic space, then $f$ is a
	smooth-Nisnevich covering if and only if $f$ admits a splitting
	cover.
\end{lemma}
\begin{proof}
	Clearly, if $f$ has a splitting cover, then it is a
	smooth-Nisnevich covering. For the other direction: let $s\in |S|$
	be a point. By Example \ref{ex:alg-sp},
	$f_s \colon X\times_S \Spec \kappa(s) \to \Spec \kappa(s)$ admits a
	section. By \cite[Lem.~2.1]{MR3754421} and \cite[Prop.~B.2 \&
	B.3]{RydhApproximation}, there exists an immersion $V_s \hookrightarrow S$
	of finite presentation such that $s\in |V_s|$ and a section to
	$f_{V_s} \colon X\times_S V_s \to V_s$. Since the $V_s$ are
	constructible, $S$ is covered by finitely many and so there are
	$s_j \in |S|$, $j=1$, $\dots$, $r$ such that
	$\cup_{j=1}^{r} |V_{s_j}| = |S|$. % By passing to quasi-compact
	% immersions inside the $V_{s_j}$, we obtain a splitting partition of
	% $S$. 
	% Now argue as in \cite[Proof of
	% Prop.~4.4]{MR2774654} to produce the filtration and a splitting
	%sequence.
\end{proof}
% \begin{lemma}
%   Let $f \colon X \to S$ be a smooth morphism of algebraic stacks.
%   \begin{enumerate}
%   \item Let $S' \to S$ be a morphism of algebraic stacks. If $f$ is
%     Nisnevich smooth covering, then $f' \colon X\times_S S'\to S'$ is a Nisnevich
%     smooth covering.
%   \item Let $T$ be a scheme and $G \to T$ a smooth and special group
%     scheme over $T$ (e.g., $G=\mathrm{GL}_{N,T}$). If $X$ is an algebraic
%     space over $T$ with an action of $G$ over $T$ such that
%     $S \simeq [X/G]$, then $f$ is a Nisnevich smooth covering.
%   \end{enumerate}
% \end{lemma}
% \begin{lemma}
% \end{lemma}
The following Proposition leverages splitting coverings to obtain a boundedness result.
\begin{proposition}\label{P:boundedness}
	Consider a cartesian square of algebraic stacks:
	\[
	\xymatrix{X' \ar[r]^{x} \ar[d]_{f'} & X \ar[d]^f \\ S' \ar[r]_s &
		S.}
	\]
	Assume that
	\begin{enumerate}
		\item $f$ and $s$ are smooth-Nisnevich coverings; and
		\item $X$ is a quasi-compact and quasi-separated algebraic space.
		%  \item $S'$ is an algebraic space. 
	\end{enumerate}
	Then there exists a quasi-compact open $S'' \subseteq S'$ such that
	$S'' \to S$ is a smooth-Nisnevich covering.
\end{proposition}
\begin{proof}
	Since $X$ is a quasi-compact algebraic space and $x$ is a
	smooth-Nisnevich covering, it follows from Lemma
	\ref{L:split-covering-sm} that $x$ admits a splitting cover
	$|X| = \cup_{j=1}^r |W_j|$ with sections
	$t_j \colon W_j \to X\times_S W_j$ to
	$f_{W_j} \colon X\times_S W_j \to W_j$. % That is, there is a finite
	% filtration of $X$ by quasi-compact open subspaces
	% $\emptyset = X_0 \subseteq \cdots \subseteq X_r = X$ and sections
	% $t_i$ to $x_{Z_i} \colon X' \times_X Z_i \to Z_i$ for each $i=1$,
	% $\dots$, $r$, where $Z_i = S_i\setminus S_{i-1}$. 
	Since $x$ is quasi-separated and locally of finite presentation, the
	sections $t_j$ are of finite presentation. But the $W_j$ are
	quasi-compact, so the subsets $t_j(W_j) \subseteq X'$ are all
	constructible (Chevalley's Theorem). Write $S' = \cup_{i\in I} S_i$
	as an increasing union of quasi-compact open subsets. Since
	$t_j(W_j)$ is constructible and $f'$ is quasi-compact, for each $j$
	there is an $i_j$ such that
	$t_j(W_j) = f^{-1}(S') \cap t_j(W_j) = f^{-1}(S_{i_j}) \cap
	t_j(W_j)$; in other words, $t_j(W_j) \subseteq
	f^{-1}(S_{i_j})$. Take $v = \max\{i_j\}$ and set $S'' = S_{v}$ and
	$X'' = f'^{-1}(S'') \subseteq X$. Then the induced morphism
	$X'' \to X$ is a smooth-Nisnevich covering (Lemma
	\ref{L:split-covering-sm}). Thus, $X'' \to S$ is a smooth-Nisnevich
	covering and so $S'' \to S$ is a smooth-Nisnevich covering.
\end{proof}

% \section{ \'Etale families in a smooth morphism}
Let $f \colon X \to S$ be a morphism of algebraic stacks. Let
$\underline{\mathrm{HS}}_f \to S$ be the \emph{Hilbert stack} of $f$,
which parameterizes quasi-finite and representable morphisms to $X$
that are proper over the base. If $f$ has affine stabilizers with
quasi-compact and separated diagonal, then
$\underline{\mathrm{HS}}_f \to S$ is a morphism of algebraic stacks
with quasi-affine diagonal, which is locally of finite presentation if
$f$ is so \cite{MR3148551,MR3359029}. As noted in
\cite[Thm.~5.1]{MR3754421}, if we let
$\underline{\mathrm{HS}}_f^{\mathrm{qfb}} \subseteq \underline{\mathrm{HS}}_f$
be the open substack parameterizing those families that are
quasi-finite over $S$, then these existence results are much simpler
to prove. Let
$\underline{\mathrm{HS}}_f^{\mathrm{\acute{e}tb}} \subseteq
\underline{\mathrm{HS}}_f^{\mathrm{qfb}}$ be the substack parameterizing those
families that are \'etale over $S$. If $f$ is a morphism of algebraic spaces, then
$\underline{\mathrm{HS}}_f^{\mathrm{\acute{e}tb}} \simeq [(X/S)^d/S_d]$
\cite[Thm.~5.1]{MR2821738}, where $(X/S)^d$ denotes the $d$-fold fiber
product of $X$ over $S$ and $S_d$ acts on this product via
permutation. If $f$ is a separated morphism of algebraic spaces, then there is an
open subspace
$\underline{\mathrm{Hilb}}_f^{\mathrm{\acute{e}tb}} \subseteq
\underline{\mathrm{HS}}_f^{\mathrm{\acute{e}tb}}$ parameterizing closed immersions
into $X$ and
$\underline{\mathrm{Hilb}}_f^{\mathrm{\acute{e}tb}} \subseteq
\underline{\mathrm{HS}}_f^{\mathrm{\acute{e}tb}}$ corresponds to the open subset of
$[(X/S)^d/S_d]$ that is the complement of the diagonals
\cite[Thm.~5.1]{MR2821738}. The following result summarizes the
relevant results of \cite{MR2821738} and can be viewed as a smooth
variant of \cite[\S\S 5-6]{MR3754421}.
\begin{proposition}\label{P:hs-covering}
	Let $f \colon X \to S$ be a representable smooth covering of algebraic stacks. 
	\begin{enumerate}
		\item \label{PI:hs-covering:HS} $\underline{\mathrm{HS}}_f^{\mathrm{\acute{e}tb}} \to S$ is
		% \begin{enumerate}
		% \item \label{PI:hs-covering:rel-dm}
		relatively Deligne--Mumford with separated diagonal, and a
		smooth-Nisnevich covering.
		\begin{enumerate}
			\item \label{PI:hs-covering:HS:et} If $f$ is \'etale, then
			$\underline{\mathrm{HS}}_f^{\mathrm{\acute{e}tb}} \to S$ is \'etale.\footnote{If
				$X$ and $S$ are quasi-separated, then $f$ is even a Nisnevich
				covering \cite[Proof of Thm.~4.1]{MR3754421}.}
			\item \label{PI:hs-covering:HS:dm} If $X$ is a Deligne--Mumford
			stack (with separated diagonal), then so is
			$\underline{\mathrm{HS}}_f^{\mathrm{\acute{e}tb}}$.
		\end{enumerate}
		\item \label{PI:hs-covering:hilb} If $f$ is separated, then
		$\underline{\mathrm{Hilb}}_f^{\mathrm{\acute{e}tb}} \to S$ is representable,
		separated, and a smooth-Nisnevich covering.
		\begin{enumerate}
			\item \label{PI:hs-covering:hilb:alg} If $X$ is a (separated) algebraic space, then so is
			$\underline{\mathrm{Hilb}}_f^{\mathrm{\acute{e}tb}}$.
			\item \label{PI:hs-covering:hilb:sch} If $X$ is a (separated) scheme and $f$ is \'etale, then so
			is $\underline{\mathrm{Hilb}}_f^{\mathrm{\acute{e}tb}}$.
		\end{enumerate}
		% \item \label{PI:hs-covering:nis}
		%   $\underline{\mathrm{HS}}_f^{\mathrm{\acute{e}tb}} \to S$ is a smooth-Nisnevich
		%   covering; and
		% \item \label{PI:hs-covering:abs}
	\end{enumerate}
	%  \end{enumerate}
\end{proposition}
\begin{proof}
	The first parts of \eqref{PI:hs-covering:HS} and
	\eqref{PI:hs-covering:hilb} are smooth-local on $S$ and follow
	trivially from the explicit descriptions above when $f$ is a
	morphism of algebraic spaces. % Hence, we may
	% assume that $S$ is an affine scheme and in case
	% \eqref{PI:hs-covering:HS} that $X$ is an algebraic space that is
	% locally of finite presentation over $S$ and in
	% \eqref{PI:hs-covering:hilb} $X$ is, in addition, separated. Now
	% \cite[Thm.~5.1]{MR2821738} says that
	% $\underline{\mathrm{HS}}_f^{\mathrm{\acute{e}tb}} \simeq [(X/S)^d/S_d]$, which is
	% certainly smooth and relatively Deligne--Mumford with separated
	% diagonal. This explicit description proves
	% \eqref{PI:hs-covering:HS:et} too. Also, \cite[Thm.~5.1]{MR2821738}
	% implies that
	% $\underline{\mathrm{Hilb}}_f^{\mathrm{\acute{e}tb}} \subseteq
	% \underline{\mathrm{HS}}_f^{\mathrm{\acute{e}tb}}$ corresponds to the open subset of
	% $[(X/S)^d/S_d]$ that is the complement of the diagonals; in
	% particular, $\underline{\mathrm{Hilb}}_f^{\mathrm{\acute{e}tb}}$ is a smooth and
	% separated algebraic space over $S$.
	%We know that
	% $\eta_f \colon \underline{\mathrm{HS}}_f^{\mathrm{\acute{e}tb}} \to S$ is locally of
	% finite presentation, so the smoothness follows from the
	% infinitesimal lifting criterion (e.g.,
	% \cite[Prop.~4.15]{MR1771927}). Hence, we must solve the following
	% lifting problem: if $A' \to A$ is a surjection of strictly henselian
	% local rings such that $I=\ker(A' \to A)$ is square 0, then complete
	% the diagram
	% \[
	%   \xymatrix{\Spec A \ar[r] \ar@{_(->}[d] &
	%     \underline{\mathrm{HS}}_f^{\mathrm{\acute{e}tb}} \ar[d] \\ \Spec A' \ar@{-->}[ur]\ar[r] & S.}
	% \]
	% That is, it suffices to solve the following deformation problem:
	% \[
	%   \xymatrix{Z \ar[d]_i \ar@{^(-->}[r]^j & Z' \ar@{-->}[d] \\
	%     X\times_S \Spec A  \ar[d]_{f_A} \ar[r] & X \times_S \Spec A' \ar[d]^{f_{A'}} \\
	%     \Spec A \ar[r] & \Spec A',}
	% \]
	% where all squares are cartesian and $Z' \to \Spec A'$ is finite
	% \'etale. The inclusion
	% $\underline{\mathrm{HS}}_f^{\mathrm{\acute{e}tb}} \subseteq
	% \underline{\mathrm{HS}}_f$ is an open immersion and $f$ is flat, so
	% by \cite[\S9]{MR3589351}, the obstruction to the solution of this
	% deformation problem belongs to the group
	% $\Ext^2_{\Orb_{Z}}(L_{i},i^*f_A^*I)$; if the obstruction vanishes,
	% then the set of lifts is a torsor under
	% $\Ext^1_{\Orb_{Z}}(L_{i},i^*f_A^*I)$; and the group of automorphisms
	% of a lift is $\Ext^0_{\Orb_{Z}}(L_{i},i^*f_A^*I)$. Now corresponding
	% to $Z \to X\times_S \Spec A \to \Spec A$, there is the distinguished
	% triangle: $i^*L_{f_A} \to L_{Z/\Spec A} \to L_i \to i^*L_{f_A}[1]$.
	% Since $Z \to \Spec A$ is \'etale, $L_{Z/\Spec A} \simeq 0$. Also,
	% $f_A$ is smooth and representable so
	% $L_{f_A} \simeq \Omega_{f_A}[0]$, which is a vector bundle of finite
	% rank. Hence, $L_i \simeq i^*\Omega_{f_A}[1]$. It follows that
	% $ \Ext^r_{\Orb_Z}(L_i,i^*f_A^*I) \simeq
	% \Ext^{r-1}_{\Orb_Z}(i^*\Omega_{f_A},i^*f_A^*I) =0$ whenever
	% $r\neq 1$. That is, the diagram admits a completion and so $\eta_f$
	% is smooth. This also shows that $\eta_f$ is relatively
	% Deligne--Mumford: we know already that $\eta_f$ has quasi-affine
	% diagonal of finite presentation and
	% $\Ext^0_{\Orb_Z}(L_i,i^*f_A^*I) = 0$ shows that the diagonal is
	% formally unramified (e.g., \cite[Cor.~6.14]{MR3589351}).
	To see that $\underline{\mathrm{HS}}_f^{\mathrm{\acute{e}tb}} \to S$ is a
	smooth-Nisnevich covering, let $s \colon \Spec K \to S$ be a
	morphism, where $K$ is a field. Since $f\colon X \to S$ is a smooth
	covering, there is a finite separable field extension
	$K \subseteq K'$ together with a lift $x \colon \Spec K' \to X$
	of $s$. That is, we have
	$\Spec K' \to X \times_S \Spec K \to \Spec K$, which corresponds
	to a lift of $x$ to $\underline{\mathrm{HS}}_f^{\mathrm{\acute{e}tb}}$. If $X \to S$
	is separated, then $\underline{\mathrm{Hilb}}_f^{\mathrm{\acute{e}tb}} \to S$ is a
	smooth-Nisnevich covering as we may replace $\Spec K'$ to be the
	residue field of the closed point in its image in
	$X\times_S \Spec K$. Next, the universal family gives us a diagram:
	\[
	\xymatrix@R-1pc{\mathcal{Z} \ar[r] \ar[dr] & X\times_S
		\underline{\mathrm{HS}}_f^{\mathrm{\acute{e}tb}} \ar[r] \ar[d] & X \ar[d] \\ &
		\underline{\mathrm{HS}}_f^{\mathrm{\acute{e}tb}} \ar[r] & S.}
	\]
	If $X$ is Deligne--Mumford (with separated diagonal), then
	$X\times_S \underline{\mathrm{HS}}_f^{\mathrm{\acute{e}tb}}$ is Deligne--Mumford
	(with separated diagonal). The map
	$\mathcal{Z} \to X\times_S \underline{\mathrm{HS}}_f^{\mathrm{\acute{e}tb}}$ is
	quasi-finite, separated and representable and so $\mathcal{Z}$ is
	Deligne--Mumford (with separated diagonal). But
	$\mathcal{Z} \to \underline{\mathrm{HS}}_f^{\mathrm{\acute{e}tb}}$ is finite \'etale
	and surjective and so $\underline{\mathrm{HS}}_f^{\mathrm{\acute{e}tb}}$ is
	Deligne--Mumford (with separated diagonal). Finally,
	\eqref{PI:hs-covering:hilb:alg} follows from
	\cite[Thm.~4.1]{MR2821738} and \eqref{PI:hs-covering:hilb:sch}
	follows from \cite[Rem.~2.3 \& Thm.~2.4]{MR2821738}.
\end{proof}
%\begin{remark}
% Proposition \ref{P:hs-covering} can be trivially generalized to the
% situation where $f$ is tame, quasi-separated, and relatively
% Deligne--Mumford with separated diagonal.
%\end{remark}
\begin{proof}[Proof of Theorem \ref{MT}]
	For the existence: let $v \colon V \to \mathcal{X}$ be a smooth cover,
	where $V$ is a scheme (take $v$ to be \'etale if $\mathcal{X}$ is
	Deligne--Mumford). Then $v$ is representable and let
	$w \colon W=\underline{\mathrm{HS}}_v^{\mathrm{\acute{e}tb}} \to \mathcal{X}$ be the
	induced morphism. Proposition
	\ref{P:hs-covering}\eqref{PI:hs-covering:HS} implies that $w$ is a
	smooth-Nisnevich covering and $W$ is a Deligne--Mumford stack with
	separated diagonal (if $\mathcal{X}$ is a Deligne--Mumford stack, $w$
	is even \'etale). Replacing $\mathcal{X}$ by $W$ we are reduced to the
	situation where $\mathcal{X}$ is Deligne--Mumford with separated
	diagonal. In this case, $v$ is separated, representable, and
	\'etale. Then Proposition
	\ref{P:hs-covering}\eqref{PI:hs-covering:hilb} implies that
	$p \colon X = \underline{\mathrm{Hilb}}^{\mathrm{\acute{e}tb}}_v \to \mathcal{X}$ is
	an \'etale smooth-Nisnevich covering and $X$ is a separated scheme.
	
	For the boundedness: by
	\cite[Prop.~2.6(i)]{MR3436239},
	$|\mathcal{X}|$ admits a finite partition
	$\amalg_{j=1}^r|\mathcal{W}_j|$, where
	$\mathcal{W}_j = [W_j/\mathrm{GL}_{n_j}]$ and $W_j$ is a quasi-affine
	scheme. For each $j$ form the following cartesian diagram:
	\[
	\xymatrix@R-1pc{W_j' \ar[r]^{p_j'} \ar[d]_{q_j'} & W_j \ar[d]^{q_j} \\
		X\times_{\mathcal{X}} \mathcal{W}_j \ar[r]_{p_j} & \mathcal{W}_j.}
	\]
	Then $p_j$ is a smooth-Nisnevich covering and so too is $q_j$
	(Example \ref{ex:gln-nis-sm}). By Proposition \ref{P:boundedness}, it
	follows that there is a quasi-compact open
	$W''_j \subseteq X\times_{\mathcal{X}} \mathcal{W}_j$, which is also a
	scheme, such that $W''_j \to \mathcal{W}_j$ is a
	smooth-Nisnevich cover. Now take a quasi-compact open
	$X'' \subseteq X$ such that
	$W''_j \subseteq X'' \times_{\mathcal{X}} \mathcal{W}_j$ for all
	$j$. Then $X'' \to \mathcal{X}$ is a smooth-Nisnevich cover, which
	proves the claim.
\end{proof}
%\noindent{\bf Acknowledgements.}









\section{Applications to Motivic homotopy theory of algebraic stack}\label{section-applications}
%We will now recall the notion of the (unstable) $\A^1$-homotopy category over a field. 
Let $k$ be a field. Let $Sm/k$ denote the category of smooth schemes over $k$. Let $\Delta^{op}PSh(Sm/k)$ be the category of simplicial presheaves on $Sm/k$. Then $\Delta^{op}PSh(Sm/k)$ has a local model structure with respect to the Nisnevich topology (see \cite{Jar}). A morphism $f:X\rightarrow Y$ in $\Delta^{op}PSh(Sm/k)$ is a weak equivalence if the induced morphisms on stalks (for the Nisnevich topology) are weak equivalences of simplicial sets. Cofibrations are monomorphisms, and fibrations are characterised by the right lifting property. 

We Bousfield localise this model structure with respect to the class of maps $X\times \A^1\rightarrow X$ (see \cite[3.2]{MV}). The resulting model structure is called the Nisnevich motivic model structure. Denote by $\mathcal{H}(k)$ the resulting homotopy category. This is the (unstable) $\A^1$-homotopy category for smooth schemes over $k$.

Let us also rapidly recall the definition of Voevodsky's triangulated category of motives $\DM{k,\Z}{}$. Let $Cor_k$ denote the category of finite correspondences whose objects are smooth separated schemes over $k$. For any two $X,Y$, the morphisms of $Cor_k$ are given by $Cor(X,Y)$ which is the free abelian group generated by irreducible closed subschemes $W\subset X\times Y$ that are finite and surjective over $X$. An additive functor $F:Cor_k^{op}\rightarrow \mathbf{Ab}$ is called a presheaf with transfers. Let $PST(k,\Z)$ denote the category of presheaves with transfers. For any smooth scheme $X$, let $\Z_{tr}(X)$ be the presheaf with tranfers which on any smooth scheme $Y$ is defined as
\[\Z_{tr}(X)(Y):=Cor(Y,X)\]
Let $K(PST(k,\Z))$ denote the category of complexes of presheaves with transfers. The category $K(PST(k,\Z))$ also has a Nisnevich motivic model structure which is defined analogously as in the case of $\Delta^{op}PSh(Sm/k)$. We denote the associated homotopy category by $\DM{k,\Z}{}$. This is Voevodsky's triangulated category of mixed motives in the Nisnevich topology (for details, see \cite{MWV}).

We now have the following trivial observation.
\begin{lemma}\label{lemma-lift-hensel-local-points}
  Let $p \colon X\rightarrow\sX$ be a smooth-Nisnevich covering of algebraic stacks. If
  $\Spec \sO$ is the spectrum of a Henselian local ring $\sO$, then
  every morphism $\Spec \sO\rightarrow \sX$ factors through $p$.
\end{lemma}
\begin{proof}
	Let $K$ be the residue field of $\sO$ and $X_\sO:= X\times_\sX \Spec \sO$ denote the base change of $X$ to $\Spec \sO$. Then the composite $\Spec K \rightarrow \Spec \sO \rightarrow \sX$ lifts to a morphism $\Spec K \rightarrow X$. This also gives us a morphism $\Spec K \rightarrow X_\sO$. Now, as $X_\sO \rightarrow \Spec \sO$ is smooth, we have a surjection of sets $X_\sO(\sO)\twoheadrightarrow X_\sO(K)$. Thus, we have a section $\Spec \sO \rightarrow X_\sO$. Composing with the projection map $X_\sO \rightarrow X$ gives us the desired lift.
\end{proof}





\begin{proposition}
	Let $\sX$ be an algebraic stack that is locally of finite type over a field $k$. Let $X\rightarrow \sX$ be a smooth-Nisnevich covering. Then the associated \v{C}ech hypercover $X_{\bullet}$ is weakly equivalent to $\sX$ in the Nisnevich local model structure on the category $\Delta^{op}PSh (Sm/k)$ of simplicial presheaves.
\end{proposition}
\begin{proof}
	By Lemma \ref{lemma-lift-hensel-local-points}, we know that hensel local points lift along smooth-Nisnevich coverings. Adapting the proof of \cite[Theorem 1.2]{cdh20} gives the result.
\end{proof}

The proof of Theorem \ref{corollary-smooth-nisnevich-presentation} follows from $\A^1$-localising the above proposition.

\begin{proof}[Proof of Theorem \ref{corollary-smooth-nisnevich-presentation}]
	As $\A^1$-localisation preserves simplicial equivalences, the result follows from the previous proposition.
\end{proof}

As a consequence, we have the following result.

\begin{corollary}
	The weak equivalence above induces an equivalence of motives $M(X_{\bullet})\simeq M(\sX)$ in $\DM{k,\Z}{}$.
\end{corollary}
\begin{proof}
	Use the functor $M: \sH_{\bullet}(k)\rightarrow \DM{k,\Z}{}$ (see \cite{cdh20} for details).
\end{proof}

Combining these results with Theorem \ref{MT}, we obtain motivic
results for smooth algebraic stacks.  The proofs are exactly as in
\cite{cdh20} after replacing $\mathrm{GL}_n$-presentations with
smooth-Nisnevich presentations.
\begin{theorem}
	Let $\sX$ be an algebraic stack that is smooth over a field $k$. Then its motive $M(\sX)\in \DM{k,\Z}{}$ satisfies the following properties:
	\begin{enumerate}
		\item $M(\sX)$ satisfies Nisnevich descent.
		\item (Projective bundle formula) For any vector bundle $\sE$ of rank $n+1$ over $\sX$, we have a canonical isomorphism
		\[M(\Proj(\sE))\simeq \bigoplus_{i=0}^n M(\sX)(i)[2i]\]
		\item (Blow-up formula) For $\sZ \subset \sX$ a smooth closed substack of pure codimension $c$ we have
		\[M(Bl_{\sZ}(\sX))\simeq M(\sX)\oplus_{i=0}^{c-1} M(\sZ)(i)[2i]\]
		\item (Gysin triangle) For $\sZ \subset \sX$ a smooth closed substack of codimension c, we have a Gysin triangle:
		\[M(\sX \setminus \sZ)\rightarrow M(\sX)\rightarrow M(\sZ)(c)[2c] \rightarrow M(\sX \setminus \sZ)[1].\]
	\end{enumerate}
\end{theorem}
\begin{proof}
	The proofs in \cite[\S 4]{cdh20} go through using a smooth-Nisnevich covering $X\rightarrow\sX$.
\end{proof}

\begin{remark}
	Note that the transfer functor $M: \sH_{\bullet}(k)\rightarrow \DM{k,\Z}{}$ allows us to define the motive of \textit{any} smooth stack. What is apriori unclear is whether this motive has any properties (good or bad). This is why in \cite{cdh20}, the authors work with stacks that are Nisnevich locally quotient stacks in order to use homotopical descent to prove the above formulae for the motive. With Theorem \ref{corollary-smooth-nisnevich-presentation}, those arguments now work for all smooth stacks.
\end{remark}


%by stating two more results about the motivic cohomology of smooth stacks. The first one states that the motivic cohomology can be computed as hypercohomology on the lisse-Nisnevich site of the stack; while the second one is a comparison theorem between the motivic cohomology and \'{e}tale cohomlogy with finite coefficients.


%\section{Motivic cohomology}
\subsection{Motivic cohomology as hypercohomology}
We briefly recall the definitions of smooth-Nisnevich and
smooth-Zariski sites of an algebraic stack.
\begin{definition}
	Let $\sX$ be an algebraic stack. The smooth-Nisnevich (resp. smooth-Zariski) site of $\sX$, denoted by $\sX_{\text{lis-nis}}$ (resp.\ $\sX_{\text{lis-zar}}$) is the category whose objects consist of pairs $(U,p)$ where $p$ is a smooth morphism $p:U\rightarrow \sX$  from an algebraic space $U$ and coverings in $\sX_{\text{lis-nis}}$ are given by Nisnevich coverings (resp. Zariski coverings) of algebraic spaces.
\end{definition}

The following result states that motivic cohomology can be computed as
hypercohomology of the motivic complexes $\Z(i)$ on the
smooth-Nisnevich or smooth-Zariski (in the Deligne--Mumford case) site
of the stack.

\begin{proposition}\label{proposition-hypercohomology}
	Let $\sX$ be an algebraic stack that is smooth over a field $k$. The motivic cohomology of $\sX$ agrees with the hypercohomology of the motivic complexes $\Z(j)$ on $\sX_{\text{lis-nis}}$. That is,
	\[Ext^i_{D(\sX)}(\Z,\Z(j)|_\sX)\simeq
          Ext^i_{DA(\sX)}(\Z,\Z(j)|_\sX)\simeq
          Hom_{\DM{k,\Z}{}}(M(\sX),\Z(j)[i]),\] where $\Z$ denotes the
        constant sheaf $\Z$ on the $\sX_{\text{lis-nis}}$. In
        addition, if $\sX$ is separated and Deligne--Mumford with
        schematic coarse space, then the motivic cohomology of $\sX$
        agrees with the hypercohomology of the motivic complexes
        $\Z(j)$ on $\sX_{\text{lis-zar}}$.
\end{proposition}
\begin{proof}
  See \cite[Theorem 5.2]{cdh20} for the smooth-Nisnevich case. In the
  Deligne--Mumford case, it follows from \cite{KreschGeometry}, where it is shown that such stacks are Zariski-locally quotient stacks.
\end{proof}

% In fact, for smooth separated Deligne-Mumford stacks with schematic coarse space over a field $k$, we can also use the smooth-Zariski site for the above computation. This follows from the following obvious corollary of \cite{KreschVistoli}:

% \begin{corollary}
% 	Let $\sX$ be a separated Deligne-Mumford stack that is smooth over a field $k$. If $\sX$ has   schematic coarse space, then $\sX$ is Zariski locally a quotient stack.
% \end{corollary}
% \begin{proof}
% 	By \cite{KreschVistoli}, it is immediate that $\sX$ admits finite flat covers Zarski locally on the coarse space. As resolution property descends along finite flat maps, we are done.
% \end{proof}

% The above propostion show that for any such stack one can find a $GL_n$-presentation which is a \emph{Zariski local} equivalence.

% We note the following refinement of Corollary \ref{theorem-hypercohomology} which is true for smooth separated Deligne-Mumford stacks with schematic coarse space. The proof is that same using a Zariki local presentation (instead of a Nisnevich local presentation).

% \begin{corollary}%\label{theorem-hypercohomology}
% 	Let $\sX$ be a smooth separated Deligne Mumford stack over a field $k$. Assume that the coarse moduli space of $\sX$ is a scheme. The motivic cohomology of $\sX$ agrees with the hypercohomology of the motivic complexes $\Z(j)$ on $\sX_{\text{lis-zar}}$. That is,
% 	\[Ext^i_{D(\sX)}(\Z,\Z(j)|_\sX)\simeq Ext^i_{DA(\sX)}(\Z,\Z(j)|_\sX)\simeq Hom_{\DM{k,\Z}{}}(M(\sX),\Z(j)[i]),\]
% 	where $\Z$ denotes the constant sheaf $\Z$ on the $\sX_{\text{lis-nis}}$.
% \end{corollary}


\subsection{Motivic cohomology with finite coefficients} We also have the following comparison theorem relating motivic cohomology with $\Z/n\Z$ to \'{e}tale cohomology with $\mu_n$-coefficients.
\begin{corollary}\label{corollary-bl}
	Let $\sX$ be an algebraic stack that is smooth over a field $k$. Then the homomorphisms
	\[H^{p,q}_M (\sX,\Z/n\Z)\rightarrow H^p_{\acute{e}t}(\sX,\mu_n^{\otimes q}),\]
	are isomorphisms for $p\leq q$ and monomorphisms for $p=q+1$.
\end{corollary}
\begin{proof}
	See \cite[Corollary 5.3]{cdh20}.
\end{proof}



\subsection{Motives with compact support}

In this subsection, we will define the notion of a motive with compact support for smooth algebraic stacks with affine stabilizers over a field $k$. The definition is very similar to Totaro's construction for quotient stacks using approximation by vector bundles \cite{TotaroChow}, except that since we cannot approximate non-quotient stacks by algebraic spaces, we are forced to work with homotopy limits. We will use use the definition of dimension for algebraic stacks as in \cite[Tag \spref{0AFL}]{stacks-project}

\begin{definition}
	Let $\sX$ be a connected algebraic stack of finite type over a field $k$ with affine stabilizers. Consider a smooth-Nisnevich covering $p:X\rightarrow\sX$, where $X$ is of finite type over $k$ (Theorem \ref{MT}) and let $n$ denote the relative dimension of $p: X\rightarrow\sX$. Let  $p_{\bullet}:X_{\bullet}\rightarrow \sX$ be the \v{C}ech nerve of $p$. Then we define the \textit{compactly supported motive} of $\sX$ as
	\[M^c(\sX):= \holim\!_i M^c(X_i)(-(i+1)n^2)[-2n^2i].\]
	%Here the limit is taken over the category of hypercovers of $\sX$.
\end{definition}

This definition is a bit weird looking, but it has the following pleasant feature, which also shows that the definition is independent of choices for smooth algebraic stacks.

\begin{theorem}[Poincar\'{e} Duality]
	Let $\sX$ be a smooth algebraic stack of dimension $d$. Then we have an isomorphism, 
	\[M^c(\sX) = M(\sX)^\vee(d)[2d].\]
\end{theorem}

\begin{proof}
	The proof is a straightforward manipulation of the definition. Tensoring by the dualizable object $\Z(-d)[-2d]$ we get:
	\begin{align*}
	M^c(\sX)\otimes \Z(-d)[-2d] &= (\holim\!_i M^c(X_i)(-(i+1)n^2)[-2n^2i])\otimes \Z(-d)[-2d]\\
	&= \holim\!_i M^c(X_i)(-(i+1)n^2-d)[-2n^2i-2d]\\
	&= \holim\!_i M(X_i)^\vee.
	\end{align*}
	where the last step follows from the fact that each $X_i$ is smooth. Further, since $M(\sX)\simeq \hocolim\!_iM(X_i)$, taking dual we get that $M(\sX)^\vee=\holim\!_i M(X_i)^\vee$. Thus, we have proved that
	\[M^c(\sX)(-d)[2d]\simeq M(\sX)^\vee.\]
	Tensoring by $\otimes \Z(d)[2d]$, gives the required expression.
\end{proof}

\noindent In the following subsection, we note an application to stable motivic homotopy theory of stacks. We will not explain any details to prevent drowning the reader in $\infty$-categories. Our intention is to simply indicate how the geometric content of Theorem \ref{MT} can be applied in the world of motivic homotopy theory. The interested reader may consult \cite{chow21} or \cite{khan2021generalized} for further details.
\subsection{The stable motivic homotopy category of a stack}
Recently, two different definitions of the stable motivic homotopy category for stacks have appeared in literature. The limit-extended category $SH_\lhd(\sX)$ in \cite{khan2021generalized} and the category $SH_{ext}^{\otimes}(\sX)$ in \cite{chow21} that is extended from schemes to stacks having Nisnevich local sections. In \cite{chow21} it is proved that the two definition are equivalent
whenever the stack admits a smooth-Nisnevich covering. As coverings such always exist for algebraic stacks by Theorem \ref{MT}, we have a strengthening of \cite[Corollary 2.5.4]{chow21}.

\begin{definition}
	Let $\sX$ be an algebraic stack. A smooth presentation $X\rightarrow\sX$ is said to admit \textit{Nisnevich local sections} if, for any scheme $Y$ and a morphism $Y\rightarrow\sX$, there exists a Nisnevich covering $Y'\rightarrow Y$ with a section $Y'\rightarrow Y\times_\sX X$. Thus, we have the following commutative diagram,
	\begin{center}
		\begin{tikzcd}
			& Y\times_\sX X\arrow[d]\arrow[r] & X\arrow[d]\\
			Y' \arrow[ur, dotted]\arrow[r] & Y\arrow[r] & \sX
		\end{tikzcd}
	\end{center}
\end{definition}

The following lemma shows that the above definition is equivalent to the notion of a smooth-Nisnevich covering.

\begin{lemma}
	Let $\sX$ be an algebraic stack. A smooth presentation $f: X\rightarrow \sX$ is admits Nisnevich local sections if and only if it is a smooth-Nisnevich covering.
\end{lemma}
\begin{proof}
	The only if part is clear. To see the if part, observe that if $X\rightarrow\sX$ is a smooth-Nisnevich covering, then by Cor 2.7, for any Henselian local ring $\sO$ with a map $\Spec \sO\rightarrow \sX$ we have a lift $\Spec\sO \rightarrow X$. Thus, if $V\rightarrow\sX$ is a morphism and $\sO_{v,V}$ is the Henselian local ring at a point $v\in V$, we have a lift $\Spec\sO_{v,V}\rightarrow V\times_\sX X$. By standard limit arguments [EGA], there exists a Nisnevich neighbourhood $V'$ of $v\in V$, and a section $V'\rightarrow V\times_\sX X$.
\end{proof}

\begin{corollary}
	Let $\sX$ be a quasi-separated algebraic stack. Then $SH_\lhd (\sX)\simeq SH_{ext}^{\otimes}(\sX)$, whenever they are defined.
\end{corollary}
\begin{proof}
	The proof in \cite[Corollary 2.5.4]{chow21} works verbatim using a smooth-Nisnevich covering $X\rightarrow \sX$.
\end{proof}


\begin{remark}
	Theorem \ref{MT} also improves \cite[Corollary 12.28]{khan2021generalized} in a similar fashion.
\end{remark}



\appendix




\begin{comment}
Note that if $\sX$ is a smooth scheme, then this definition agrees with the usual defintion of compactly supported motive.

\begin{theorem}
Let $X$ be a smooth scheme. Then
\[M^c(X)=\underset{X_{\bullet}}{\holim}\holim\!_i M^c(X_i)\]
\end{theorem}


\begin{theorem}[Localisation sequence]
	Let $\sX$ be a smooth algebraic stack. For any open substack $\sU\subset \sX$, we have an exact triangle
	\[M^c(\sX \setminus \sU)\rightarrow M^c(\sX)\rightarrow M^c(\sU)\]
\end{theorem}
\begin{proof}
	This follows 
\end{proof}

\end{comment}


\begin{comment}
\begin{abstract}
We construct the motive of an algebraic stack in the Nisnevich topology. For stacks which are Nisnevich locally quotient stacks, we give a presentation of the motive in terms of simplicial schemes. We also show that the motivic cohomology agrees with the Chow groups of Edidin-Graham-Totaro with integer coefficients.
\end{abstract}

\maketitle

\section{Introduction}
For a smooth scheme $X$ over a perfect field $k$, we have natural isomorphisms,
\begin{equation}
\label{equation-chow-comparison}
H^{n,i}(X,\Z)\cong CH^i(X,2i-n)
\end{equation}
between the motivic cohomology groups \cite{MWV} on the left, and Bloch's higher Chow groups \cite{Bloch86} on the right. When $n=2i$, this relates motivic cohomology groups $H^{2i,i}(X,\Z)$ to the ordinary Chow groups $CH^i(X,0)=CH^i(X)$ of $X$. %, i.e, $H^{n,i}(X,\Z)=Hom_{\DM{k,\Z}{}}(M(X), \Z(i)[n])$. 
Such a comparison theorem for ordinary Chow groups was first proved by Voevodsky in \cite{FSV} under the assumption that $k$ admits resolution of singularities. It was extended to all higher Chow groups for $k$ perfect, again by Voevodsky, in \cite{MWV}. %This theorem uses the Nisnevich topology in a crucial way as the triangulated category of motives is defined using Nisnevich sheaves with transfer.
The motivic cohomology groups can be thought of as cohomology groups of the motive $M(X)$ of $X$ in the triangulated category of motives, $\DM{k,\Z}{}$ (see Section \ref{section-motive-construction}).

The significance of the isomorphisms in (\ref{equation-chow-comparison}) is that we can translate various statements about Chow groups and intersection theory into statements about motivic cohomology. This is useful because statements and constructions involving (higher) Chow groups are often hard to prove (for example, Bloch's localisation sequence). However, arguments are greatly simplified if we consider analogous statements about motivic cohomology in $\DM{k,\Z}{}$. The two can then be related by Equation (\ref{equation-chow-comparison}).

Guided by this observation, in this article, we study the motives of algebraic stacks and derive various results about them. We also relate these motives to the Chow groups of algebraic stacks.

For algebraic stacks, the question of the correct notion of Chow groups is a subtle one. The first definition of Chow groups for algebraic stacks was given by Vistoli in \cite{VistoliIntersection} developing a rational intersection theory for stacks in the style of the Fulton-MacPherson theory for schemes. However, as there are not many invariant cycles on a stack, these Chow groups turn out to be quite small. Moreover, they lack expected properties like homotopy invariance, intersection product, etc. A more refined notion, as suggested by Totaro, was defined by Edidin and Graham in \cite{EG} for quotient stacks (and later generalised by Kresch in \cite{KreschChow} for all Artin stacks) which includes ``more cycles" and has better properties. These are called the Edidin-Graham-Totaro Chow groups and they are defined over integers. With $\Q$-coefficients, these two Chow groups agree. In this article, we relate the Edidin-Graham-Totaro Chow groups to the motivic cohomology of algebraic stacks (see Theorem \ref{chow-comparison}).

The story on the motivic side for algebraic stacks begins with Bertrand To\"{e}n. The notion of a motive for stacks was first defined by To\"{e}n for Deligne-Mumford stacks in \cite{Toen}. The general theory has been subsequently developed and extended by various authors (\!\!\cite{MotivesDM},\cite{HoyoisSixOperations},\cite{Joshua1,Joshua2}).
%The general theory for motives of algebraic stacks has been defined and studied in Joshua \cite{Joshua1,Joshua2} and Choudhury \cite{MotivesDM} (see also, \cite{TotaroChow, hoskins2019} for some explicit computations). 
However, most of the existing literature on motives for algebraic stacks deals with $\Q$-coefficients and in the \'{e}tale topology. In \cite{Joshua2}, a comparison theorem similar to Theorem \ref{chow-comparison} is proved for the \textit{\'{e}tale} motivic cohomology and Totaro's Chow groups \textit{with $\Q$-coefficents}.  We use the notation $\DM{k,\Q}{\acute{e}t}$ for the corresponding triangulated category of \'{e}tale motives.

The Nisnevich topology is better suited for the study of Chow groups from the motivic perspective. This is evident from the comparison isomorphisms (\ref{equation-chow-comparison}) of Voevodsky. But as Nisnevich and \'{e}tale motivic cohomology agree over $\Q$, the choice of topology does not present any serious difficulty for comparing motivic cohomology with Chow groups rationally. However, all torsion information in the Chow groups is lost after tensoring with $\Q$. Hence, having a theory of motives in the Nisnevich topology with integral coefficients is desirable.

One definition for such a motive has already been given in \cite{hoskins2019} by Hoskins and Lehalleur. However, their definition employs the ``finite dimensional approximation" technique of Totaro in \cite{TotaroChow}, and is restricted to what they call \textit{exhaustive stacks}. This definition is applicable for quotient stacks as well as stacks like the moduli of vector bundles on smooth curves, but fails to hold in general.
%With integral coefficients, however, the Nisnevich and \'{e}tale motivic cohomology are different. 


%From the comparison theorems of Voevodsky it is clear that Nisnevich topology is the correct setting to study Chow groups from the motivic perspective.  However, all torsion information in the Chow groups is lost after tensoring with $\Q$. Hence, having a theory of motives with integral coefficients is desirable. With integral coefficients, however, the Nisnevich and \'{e}tale motivic cohomology are different. 

The construction of a motive for an algebraic stack described in \cite{MotivesDM} (for the \'{e}tale topology) can also be carried out in the Nisnevich topology. This gives us a canonical definition of the motive for any algebraic stack locally of finite type over a field $k$. In this paper, we explore this construction. Note that while the construction defines the motive of an algebraic stack in great generality, it is ill-suited for the purpose of computations. Hence, we restrict ourselves to algebraic stacks that are ``Nisnevich locally quotient stacks", in the following sense:

\begin{definition}
	\label{definition-cd-quotient-stack}
	Let $\sX$ be an algebraic stack locally of finite type over a field $k$. We say that $\sX$ is a \textit{global quotient stack} if $\sX= [X/GL_n]$ for an algebraic space $X$. We say that $\sX$ is a \textit{cd-quotient stack} if it admits a Nisnevich covering $[X/GL_n]\rightarrow \sX$ by a global quotient stack\footnote{This nomenclature is inspired by \cite[Definition 2.1]{RydhApproximation}. Cd stands for \textit{completely decomposable}. Cd-topology is the old name for Nisnevich topology.}.
\end{definition}

For us, a Nisnevich covering of an algebraic stack will always mean a representable morphism (representable by \textit{algebraic spaces}) whose base change to any scheme is a Nisnevich covering. Note that we do not assume that the morphism is schematic (representable by schemes). Recall that a \emph{Nisnevich covering} $f: X\rightarrow Y$ of algebraic spaces is a surjective \'{e}tale morphism such that for any field-valued point $y:\Spec k\rightarrow Y$ there exists a lift $x: \Spec k\rightarrow X$ such that $y=f\circ x$.

The property of being a cd-quotient is enjoyed by a large class of algebraic stacks. Every global quotient stack is a cd-quotient stack. Further, for any quotient stack $[X/G]$ with $G$ a linear algebraic group, $X\times^G GL_n \rightarrow [X/G]$ is a $GL_n$-torsor, realising $[X/G]$ as a global quotient stack. 

Also, exhaustive stacks of \cite{hoskins2019} turn out to be cd-quotient stacks (see Section \ref{section-exhaustive}).


If $k$ is a perfect field then any stack with quasi-finite diagonal over $k$ is a cd-quotient stack (see \cite[\S 2]{conrad}). In particular, this includes all quasi-separated Deligne-Mumford stacks over $k$. 

Furthermore, if $\sX$ admits a good moduli space (in the sense of \cite{goodmodulispaces}), then by a theorem of Alper, Hall and Rydh, it is Nisnevich locally of the form $[\Spec(A)/GL_n]$ (see \cite[Theorem 13.1]{AHR19}), and hence a cd-quotient stack.


The property of being a quotient stack Nisnevich locally gives us a greater handle on the stack from the computational point of view. In particular, we can construct a presentation of the stack in terms of certain simplicial schemes in the (unstable) $\A^1$-homotopy category $\sH_{\bullet}(k)$. More precisely, we have the following:

\begin{theorem}\label{motive-construction}
	Let $\sY\rightarrow \sX$ be a representable Nisnevich cover of algebraic stacks over a field $k$. Assume that $\sY$ is of the form $[Y/GL_n]$ for some algebraic space $Y$, and let $p:Y\rightarrow\sY\rightarrow\sX$ be the composite. Then, for the \v{C}ech hypercover $Y_{\bullet}$ associated to $p $, $p_{\bullet}:Y_{\bullet}\rightarrow \sX$ is a Nisnevich local weak equivalence, i.e, $Y_{\bullet}\simeq \sX$ in $\sH_{\bullet}(k)$.
\end{theorem}

The above theorem gives us a nice description of the motive of a cd-quotient stack in terms of the motive of the simplicial scheme $Y_{\bullet}$ in $\DM{k,\Z}{}$. This allows us to reduce many computations about motives to stacks to motives of (simplicial) schemes. In fact, we will use this to show that for quotient stacks, the cohomology groups of this motive agree with the Chow groups of Edidin-Graham-Totaro (see \cite{EG}, \cite{KreschChow}). 

\begin{theorem}
	\label{chow-comparison}
	Let $\sX:=[X/GL_r]$ be a quotient stack, where $X$ is a smooth quasi-projective scheme with a smooth action of $GL_r$. Then the Edidin-Graham-Totaro (higher) Chow groups and the motivic cohomology groups agree integrally, i.e, 
	\[CH^i(\sX, 2i-n)\simeq H^{n,i}(\sX,\Z).\]
\end{theorem}

The proof is quite straightforward and follows easily from \cite[Proposition 3.2]{krishna2012} and Theorem \ref{motive-construction}. The quasi-projective hypothesis is required to ensure that quotients of schemes by $GL_n$-actions remain schemes (see \cite[\S 2.2.1]{krishna2012}).\\

%Fix a base field $k$.

%\noindent\underline{Hypothesis on Stacks:} Let $\sX$ be a stack such that $\sX$ admits a Nisnevich covering by quotient stack of the form $[X/GL_n]$\fixme{possible name: cd-quotient stack}.\\

\noindent\textbf{Outline.} The paper is structured as follows: In Section \ref{section-motive-construction}, we define the Nisnevich motive for an algebraic stack following \cite{MotivesDM}. We then prove Theorem \ref{motive-construction}. The argument involves two key ideas: that every principal $GL_n$-bundle over a Henselian local ring is trivial, and that fibre products of stacks are a model for their homotopy fibre products in the model category of presheaves of groupoids (see \cite[Remark 2.3]{Holl}). Theorem \ref{motive-construction} gives us a computational handle on cd-quotient stacks which is used in the sequel to prove various results.

We then show, in Section \ref{section-chow-comparison}, that for quotients of quasi-projective schemes by $GL_n$, the integral motivic cohomology of an algebraic stack agrees with the Edidin-Graham-Totaro Chow groups (Theorem \ref{chow-comparison}). As stated earlier, such a result has already been proven in \cite[Theorem 3.5]{Joshua2} for the \'{e}tale motivic cohomology groups. The crucial difference being that in \cite{Joshua2} the comparison theorem is proved for rational Chow groups (and for the \'{e}tale topology) whereas we prove it for integral Chow groups (and in the Nisnevich topology). So, in a sense, the additional information we get from Theorem \ref{chow-comparison} is that the motivic cohomology groups and the Edidin-Graham-Totaro Chow groups also agree in their torsion parts.

In Section \ref{section-various-triangles}, we convince ourselves that various exact triangle for motives continue to hold for our construction. In particular, we establish Nisnevich descent (Proposition \ref{nisnevich-descent}), projective bundle formula (Theorem \ref{projective-bundle-formula}) and the Gysin triangle (Theorem \ref{gysin-triangle}) for cd-quotient stacks. This section adapts the corresponding results for the \'{e}tale topology in \cite{MotivesDM}.

In Section \ref{section-exhaustive}, we compare our construction of the motive with the one in \cite{hoskins2019} for exhaustive stacks. We show that these two agree in $\DM{k,\Z}{}$. A similar comparison is proved for \'{e}tale motives in \cite[Appendix A]{hoskins2019}.

We include a technical result about model categories (Lemma \ref{homotopy-colimits-left-adjoints}) that is used for proving the projective bundle formula (Theorem \ref{projective-bundle-formula}) in an appendix. This is an elementary result which is surely known to anybody familiar with model categories, and we only include it for our own edification.\\

\noindent\textbf{Acknowledgements.} The second-named author was supported by the INSPIRE fellowship of the Department of Science and Technology, Govt.\ of India during the course of this work.


\section{Nisnevich motive of an algebraic stack}
\label{section-motive-construction}
In this section we will define the motive of an algebraic stack as well as prove Theorem \ref{motive-construction}. In order to do this, the following observation will be crucial. %The reason why these two things are presented together is Remark \ref{remark-gln-presentation}.


\begin{remark}[Stacks as simplicial sheaves]
	\label{remark-stacks-as-simplicial-sheaves}
	Given an algebraic stack $\sX$ (or, more generally, a presheaf of groupoids), one associates a simplicial sheaf to $\sX$ as follows: $\sX$ defines a (strict) sheaf of groupoids $\sX \in Fun((Sch/k)^{op}, Grpds)$, which sends any $k$-scheme $U\mapsto \sX_U$ in the category of groupoids. Applying the nerve functor objectwise, we get a sheaf of simplicial sets. Let us briefly recall this procedure.\\
	Given a groupoid $\sX_U$ its nerve is a simplicial set $N(\sX_U)$ whose $k$-simplices are given by $k$-tuples of composable arrows,
	\[N(\sX_U)_k =\{A_0\overset{f_1}{\rightarrow} \ldots \overset{f_k}{\rightarrow} A_k \; | \; \text{$A_i$'s are objects and $f_i$ are morphisms in $\sX_U$}\}\]
	Note that $0$-simplices are just objects of $\sX_U$, while $1$-simplices are morphisms between them.\\
	The face maps $d_i: N(\sX_U)_k \rightarrow N(\sX_U)_{k-1}$ are given by composition of morphism at the $i$-th object (or deleting the $i$-th object for $i=0,k$). Similarly, the degeneracy maps $s_i: N(\sX_U)_k \rightarrow N(\sX_U)_{k+1}$ are given by inserting the identity morphism at the $i$-th object. Thus, we have a functor 
	\[Fun((Sch/k)^{op}, Grpds)\rightarrow \Delta^{op} PSh(Sm/k)\]
	from the category of presheaves of groupoids to the category of simplicial presheaves	(see \cite{HollPhD} for more details). \\
	If $\sX\rightarrow \sZ$, $\sY\rightarrow \sZ$ are two morphisms of algebraic stacks, then viewing them as simplicial sheaves, one can form their homotopy fibre product $\sX \times^h_{\sZ} \sY$ in the homotopy category of simplicial presheaves. By \cite[Remark 2.3]{Holl}, the usual fibre product in the category of stacks $\sX \times_{\sZ} \sY$ serves as a model for this homotopy fibre product.
	
	In the remaining article, we will abuse notation by denoting the simplicial sheaf associated to a stack $\sX$ by $\sX$ itself.
	
\end{remark}

We will now recall the notion of the (unstable) $\A^1$-homotopy category over a field, as well as the construction of the motive of an algebraic stack as in \cite[Section 2]{MotivesDM}. 

Fix a base field $k$. Let $Sm/k$ denote the category of smooth schemes over $k$. Let $\Delta^{op}PSh(Sm/k)$ be the category of simplicial presheaves on $Sm/k$. Note that by Remark \ref{remark-stacks-as-simplicial-sheaves}, any stack $\sX$ can be considered as an object of $\Delta^{op}PSh(Sm/k)$.

$\Delta^{op}PSh(Sm/k)$ has a local model structure with respect to the Nisnevich topology (see \cite{Jar}). A morphism $f:X\rightarrow Y$ in $\Delta^{op}PSh(Sm/k)$ is a weak equivalence if the induced morphisms on stalks (for the Nisnevich topology) are weak equivalences of simplicial sets. Cofibrations are monomorphisms, and fibrations are characterised by the right lifting property. 

We Bousfield localise this model structure with respect to the class of maps $X\times \A^1\rightarrow X$ (see \cite[3.2]{MV}). The resulting model structure is called the Nisnevich motivic model structure. Denote by $\mathcal{H}_{\bullet}(k)$ the resulting homotopy category. This is the (unstable) $\A^1$-homotopy category for smooth schemes over $k$.



%The derived functor $M$ now gives us the motive for an algebraic stack.

We will now define the Nisnevich motive of an algebraic stack. We begin by rapidly recalling the construction of the triangulated category of mixed motives.

Let $Cor_k$ denote the category of finite correspondences whose objects are smooth separated schemes over $k$. For any two $X,Y$, the morphisms of $Cor_k$ are given by $Cor(X,Y)$ which is the free abelian group generated by irreducible closed subschemes $W\subset X\times Y$ that are finite and surjective over $X$. An additive functor $F:Cor_k^{op}\rightarrow \mathbf{Ab}$ is called a presheaf with transfers. Let $PST(k,\Z)$ denote the category presheaves with transfers. For any smooth scheme $X$, let $\Z_{tr}(X)$ be the presheaf with tranfers which on any smooth scheme $Y$ is defined as
\[\Z_{tr}(X)(Y):=Cor(X,Y)\]
Let $K(PST(k,\Z))$ denote the category of complexes of presheaves with transfers. The category $K(PST(k,\Z))$ also has a Nisnevich motivic model structure which is defined analogously as in the case of $\Delta^{op}PSh(Sm/k)$. We denote the associated homotopy category by $\DM{k,\Z}{}$. This is Voevodsky's triangulated category of mixed motives in the Nisnevich topology (for details, see \cite{MWV}).  Then we have a functor
\[N\Z_{tr}(-): \Delta^{op}PSh(Sm/k)\rightarrow K(PST(k,\Z))\]
which sends a simplicial scheme $X_{\bullet}$ to its normalised chain complex $N\Z_{tr}(X_{\bullet})$. The $i$-th degree term of the chain complex $N\Z_{tr}(X_{\bullet})$ is given by $\Z_{tr}(X_i)$.

Since every simplicial presheaf is weakly equivalent to a simplicial scheme (see \cite{DHI}), this determines the derived functor of $N\Z_{tr}(-)$ completely. Note that $N\Z_{tr}(-)$ is a left Quillen functor and so it admits a left derived functor $M$ on the homotopy categories:
\[M:\mathcal{H}_{\bullet}(k)\rightarrow \DM{k,\Z}{}.\]
For a scheme $X$, $M(X)$ is the image of $\Z_{tr}(X)$ in $\DM{k,\Z}{}$.

\begin{remark}
	\label{remark-motive-commutes-with-homotopy-colimits}
	$N\Z_{tr}$ is a left Quillen functor with respect to the Nisnevich motivic model structures on $\mathcal{H}_{\bullet}(k)$ and $\DM{k,\Z}{}$. Thus, Lemma \ref{homotopy-colimits-left-adjoints} implies that $M$ commutes with homotopy colimits.
\end{remark}

\begin{definition}[Motive of an algebraic stack]
	\label{definition-motive-stack}
	Let $\sX$ be an algebraic stack over $k$ thought of as a simplicial presheaf by the nerve construction outlined in Remark \ref{remark-stacks-as-simplicial-sheaves}. The motive of an algebraic stack $\sX$ is defined to be the image $M(\sX)$ in $\DM{k,\Z}{}$.\\
	Note that when $\sX$ is representable by a scheme $X$, $M(X)$ is the image of $\Z_{tr}(X)$ in $\DM{k,\Z}{}$.
\end{definition}


\begin{remark}
	The above construction was first done in \cite{MotivesDM} for the \'{e}tale model structure. However, it still goes through if we use the Nisnevich model structure instead. See also \cite{Joshua1} for an alternative approach.
\end{remark}


We will now show that cd-quotient stacks admit presentations by simplicial schemes in $\sH_{\bullet}(k)$. This is the content of Theorem \ref{motive-construction}. Having such a presentation will allow us to use homotopical descent techniques in the sequel in order to reduce various problems to the case of (simplicial) schemes. 

For the sake of clarity we will first prove Theorem \ref{motive-construction} in the case when $\sY = \sX$, i.e, when $\sX$ is a global quotient stack. Theorem \ref{motive-construction} is a minor extension of this case.

\begin{lemma}
	\label{lemma-principal-bundles-nisnevich-sections}
	Let $X$ be an algebraic space with an action of $GL_n$ and $\sX:=[X/GL_n]$ be the corresponding quotient stack. Let $X_{\bullet}$ denote the \v{C}ech hypercover associated to $p:X\rightarrow\sX$. Then the map of simplicial presheaves $p_{\bullet}: X_{\bullet}\rightarrow \sX$ is a Nisnevich local weak equivalence.
\end{lemma}
\begin{proof}	
	It suffices to check that given a hensel local ring $\sO$, the induced map on $\sO$-points $p_{\bullet}: X_{\bullet}(\Spec\sO)\rightarrow \sX(\Spec\sO)$ is a weak equivalence of simplicial sets. Further, as a stack is a $1$-truncated simplicial set (being a groupoid valued functor), $\pi_i = 0$ for $i\geq 2$. Thus, we only need to verify that $p$ induces an isomorphism of homotopy groups for $i=0,1$.
	
	\noindent\underline{$i=0$:} Any map $\Spec\sO\rightarrow \sX$ gives rise to a $GL_n$-torsor $p':X\times_{\sX} \Spec\sO \rightarrow \Spec\sO$. As $\sO$ is a Hensel local ring, any $GL_n$-torsor over it is trivial and hence admits a section. This implies surjectivity on $\pi_0$.\\
	For injectivity, let $f_1, f_2:\Spec\sO\rightarrow X$ be two $\sO$-points of $X$ such that $p(f_1)=p(f_2)$, then there exists a map $F:\Spec\sO\rightarrow X\times_{\sX}X$ which after composing with each of the projection maps becomes $f_1$ and $f_2$, respectively. The map $F$ may be thought of as a 1-simplex in the simplicial set $X_{\bullet}(\Spec\sO)$, and therefore, corresponds to a map $\Delta^1\rightarrow X_{\bullet}(\Spec\sO)$, which gives a homotopy between the points $f_1$ and $f_2$ implying injectivity.
	
	\noindent\underline{$i=1$:} In this case, we need to show that for any $\Spec\sO$-valued point of $\sX$, the homotopy fibre product with $p_{\bullet}$ is contractible. Then, the long exact sequence of homotopy sheaves gives the required isomorphism. By \cite[Remark 2.3]{Holl}, the homotopy fibre product is precisely the fibre product in the category of stacks. This is a (\v{C}ech nerve of a) trivial $GL_n$-torsor over $\Spec\sO$, i.e, given a point $\Spec\sO\rightarrow\sX$, the homotopy fibre of $p_{\bullet}$ is precisely the \v{C}ech hypercover $X_{\bullet}\times_{\sX}\Spec\sO$ corresponding to the $GL_n$-torsor $p':X\times_{\sX}\Spec\sO\rightarrow\Spec\sO$. Since $p'$ admits a section, the augmentation map $X_{\bullet}\times_{\sX}\Spec\sO\rightarrow\Spec\sO$ is a Nisnevich local weak equivalence. As, $\pi_i(\Spec\sO)=0$ for $i> 0$, we get the desired result. %\fixme{Write down somewhere why stacky fibre prodcut and homotopy fibre product agree (Maybe in an appendix).}.
\end{proof}

Any Nisnevich cover $\sY\rightarrow\sX$ admits sections Nisnevich locally. Theorem \ref{motive-construction} now follows easily from this fact and Lemma \ref{lemma-principal-bundles-nisnevich-sections}.

\begin{proof}[Proof of Theorem \ref{motive-construction}]
	We need to check that $p_{\bullet}:Y_{\bullet}\rightarrow \sX$ induces an isomorphism on all homotopy sheaves, $\pi_i$. Further, it suffices to check this on all Hensel local schemes.
	
	\noindent\underline{$i=0$:} Let $\Spec\sO\rightarrow \sX$ be a point of $\sX$. Base changing, we get maps $Y\times_{\sX}\Spec\sO\rightarrow\sY\times_{\sX}\Spec\sO\rightarrow\Spec\sO$, where the first map is a principal $GL_n$-bundle and the second is a Nisnevich cover. So, the second map admits a Nisnevich local section. By the previous lemma, so does the first. This proves surjectivity.\\
	For injectivity, let $f_1,f_2:\Spec\sO\rightarrow Y$ be two points which map to the same point in $\sX$. This implies that there exists a section $\Spec\sO\rightarrow Y\times_{\sX}Y$ which after composing with each of the projection maps becomes $f_1$ and $f_2$, respectively.
	
	\noindent\underline{$i>0$:} To show this we need to show that for any $\Spec\sO$-valued point of $\sX$, the homotopy fibre product with $p_{\bullet}$ is contractible. Then, the long exact sequence of homotopy sheaves gives us the required isomorphism.\\
	As noted in the previous lemma, the homotopy fibre product is equal to the stacky fibre product\cite[Remark 2.3]{Holl}. Hence, for a point $\Spec\sO\rightarrow\sX$, the homotopy fibre of $p_{\bullet}$ is precisely the \v{C}ech hypercover $Y_{\bullet}\times_{\sX}\Spec\sO$ of $\Spec\sO$. Since it admits a section, the augmentation map $Y_{\bullet}\times_{\sX}\Spec\sO\rightarrow\Spec\sO$ is a Nisnevich local weak equivalence. As, $\pi_i(\Spec\sO)=0$ for $i> 0$, we get the desired result.
\end{proof}




\begin{remark}
	\label{remark-gln-presentation}
	For a cd-quotient stack, we have a \v{C}ech hypercover $Y_{\bullet}\rightarrow \sX$ which is a Nisnevich local weak equivalence, by Theorem \ref{motive-construction}. Applying the functor $M:\mathcal{H}_{\bullet}(k)\rightarrow \DM{k,\Z}{}$, we see that the motive of a cd-quotient stack $\sX$ is given by the normalised chain complex $N\Z_{tr}(Y_{\bullet})$.
\end{remark}

\begin{definition}
	For a cd-quotient stack $\sX$, let $p:X\rightarrow\sX$ be the presentation obtained from the composition $X\rightarrow [X/GL_n]\rightarrow\sX$ (as in the hypothesis of Theorem \ref{motive-construction}), and let $p_{\bullet}: X_{\bullet}\rightarrow \sX$ denote the associated \v{C}ech hypercover. Motivated by the content of Theorem \ref{motive-construction}, we will call $p_{\bullet}: X_{\bullet}\rightarrow \sX$ a \textit{$GL_n$-presentation} of $\sX$. 
	%By abuse of notation, we may sometimes refer to the $0$-skeleton $p:X\rightarrow\sX$ as a motivic presentation.
\end{definition}




\begin{remark}\label{motive-not-geometric}
	The category generated by motives of stacks as constructed above is larger than the category of geometric motives (motives generated by smooth quasi-projective schemes). In fact, if $G$ is a finite group, $BG$ is not a geometric motive. To see this note that any realization of a geometric motive must have bounded cohomology. Further, for a finite group, the cohomology of $BG$ is the same as the group cohomology of the group $G$. But if $G$ is finite cyclic, then the latter is periodic in odd degrees, showing that $BG$ does not have bounded cohomology for a cyclic group.
\end{remark}

\begin{proposition}
	\label{proposition-comparison-etale-motive}
	Let $\sX$ be a cd-quotient stack. The Nisnevich motive $M(Y_{\bullet})$ of $\sX$ agrees with its \'{e}tale motive in $\DM{k,\Z}{\acute{e}t}$ after \'{e}tale sheafification.
\end{proposition}
\begin{proof}
	Let $Y_{\bullet}\rightarrow \sX$ be a $GL_n$-presentation so that $M(Y_{\bullet})\simeq M(\sX)$ in $\DM{k,\Z}{}$. Since sheafification is an exact functor, we have a functor $\sigma: \DM{k,\Z}{}\rightarrow \DM{k,\Z}{\acute{e}t}$ which takes Nisnevich local weak equivalences to \'{e}tale local weak equivalences (see \cite[Remark 14.3]{MWV}). Thus, the Nisnevich local weak equivalence $M(Y_{\bullet})\rightarrow M(\sX)$ becomes an \'{e}tale local weak equivalence after \'{e}tale sheafification. Hence, $M(Y_{\bullet})\simeq M(\sX)$ in $\DM{k,\Z}{\acute{e}t}$.
	
	In fact, more is true. We can show that any $GL_n$-presentation is equivalent to the \v{C}ech hypercovers associated to a smooth presentation of $\sX$. 
	
	Let $p: U\rightarrow \sX$ be a smooth presentation, and $p_{\bullet}: U_{\bullet}\rightarrow \sX$ be the associated \v{C}ech hypercover. Given any strict Hensel local point $\Spec \sO\rightarrow \sX$ consider the base change $p_\sO : U_{\sO}\rightarrow \Spec \sO$. This is a smooth morphism and so admits \'{e}tale local sections. In fact, since $\sO$ is a Hensel local ring, we have a section $\Spec \sO \rightarrow U_{\sO}$ of $p_\sO$. This implies that the induced map of simplicial sets $p_{\bullet}(\sO): U_{\bullet}(\sO)\rightarrow \sX(\sO)$ is a weak equivalence of simplicial sets (proof is similar to Lemma \ref{lemma-principal-bundles-nisnevich-sections}). Thus, $U_{\bullet}\rightarrow \sX$ is an \'{e}tale local weak equivalence. So the induced map on \'{e}tale motives $M(U_{\bullet}) \rightarrow M(\sX)$ is also an \'{e}tale local weak equivalence in $\DM{k, \Z}{\acute{e}t}$ (see \cite[Corollary 2.14]{MotivesDM}). Hence, $M(Y_{\bullet})\simeq M(U_{\bullet})$ in $\DM{k,\Z}{\acute{e}t}$.
	%We use the subscript $\Q$ to remember that we are dealing with $\Q$-coefficients.\\	
\end{proof}

\begin{remark}
In Proposition \ref{proposition-comparison-etale-motive}, the $GL_n$-presentation $Y_{\bullet}$ and the \v{C}ech hypercover $U_{\bullet}$ are simplicial objects in the category of \textit{algebraic spaces}. This is because a smooth presentation $p:U\rightarrow \sX$ of an algebraic stack need not be representable by schemes, but only algebraic spaces. However, as any algebraic space admits a Nisnevich presentation by a scheme \cite[Theorem II.6.4]{knutson}, we can refine the \v{C}ech hypercover $U_{\bullet}$ to a generalised hypercovering $V_{\bullet}$ such that each $V_i$ is a scheme. Then $M(V_{\bullet})$ computes the motive $M(\sX)$ (see \cite{DHI} for details).
%Note also that if the diagonal of $\sX$ is schematic (for example, if $\sX$ is a Deligne-Mumford stack) then $U_{\bullet}$ and $Y_{\bullet}$ in Proposition \ref{proposition-comparison-etale-motive} can be taken to be simplicial schemes. Then their motives $M(U_{\bullet})$ and $M(Y_{\bullet})$ are given by the normalised chain complexes $N\Z_{tr}(U_{\bullet})$ and $N\Z_{tr}(Y_{\bullet})$, respectively.
\end{remark}


\begin{remark}[\'{E}tale motives with $\Q$-coefficients]
	Tensoring with $\Q$ gives us a functor $-\otimes \Q: \DM{k,\Z}{\acute{e}t}\rightarrow \DM{k,\Q}{\acute{e}t}$ which is just change of coefficients (this also works in the Nisnevich topology). We write $M(Y_{\bullet})\otimes \Q := M(Y_{\bullet})_{\Q}$. By \cite[Theorem 14.30]{MWV}, the \'{e}tale sheafification functor $\sigma: \DM{k,\Q}{}\rightarrow \DM{k,\Q}{\acute{e}t}$ is an equivalence of categories.\\
	By \cite[Theorem 4.6]{MotivesDM}, for a smooth separated Deligne-Mumford stack $\sX$, $M(\sX)_{\Q}$ is a geometric motive. This does not contradict Remark \ref{motive-not-geometric}, since the cohomology groups of a cyclic group are torsion and will vanish after tensoring with $\Q$.\\
	Further, if $\pi:\sX\rightarrow X$ is the coarse space map, then $M(\pi)_\Q: M(\sX)_\Q\rightarrow M(X)_\Q$ is an isomorphism by \cite[Theorem 3.3]{MotivesDM}. This is clearly false integrally, since for a finite group $G$ over a field $k$, the structure map $BG\rightarrow \Spec k$ is a coarse space map.
\end{remark}


\section{Comparison with Edidin-Graham-Totaro Chow Groups}
\label{section-chow-comparison}

We will now proceed to show that the motivic cohomology groups of the motive defined by Theorem \ref{motive-construction} agree with the (higher) Chow groups defined by Edidin-Graham-Totaro for quotients of smooth quasi-projective schemes in \cite{EG}. This implicitly follows from \cite{krishna2012}, but we write out the details for the sake of completeness.
In what follows we only consider the action of $G:=GL_r$ on \textit{quasi-projective schemes}.

We begin by recalling Totaro's definition of Chow groups for quotient stacks (see \cite[Section 2.2]{EG}). The definition is via a ``Borel type construction". 

Let $X$ be an $n$-dimensional smooth quasi-projective scheme with an action of $GL_r$ and let $[X/GL_r]$ be quotient of this action. Choose an $l$-dimensional representation $V$ of $GL_r$ such that $V$ has an open subset $U$ on which $GL_r$ acts freely and whose complement has codimension greater than $n-i$. Then, we have a principal $GL_r$-bundle $X\times U\rightarrow (X\times U)/GL_r$, and we denote the quotient as $X_{GL_r}:= (X\times U)/GL_r$. Note that $X_{GL_r}$ is a quasi-projective scheme (see \cite[Lemma 2.2]{krishna2012} or \cite[Proposition 23]{EG}). 
\begin{definition}
\label{totaro-chow}
With set-up as above, we define the $i$-th Chow group of $[X/GL_r]$ as 
\[CH^i([X/GL_r]) :=CH^{i}(X_{GL_r}).\]
This definition is independent of the choice of $V$ and $U$ so long as $\codim(V,U)> i$ (see \cite[section 2.2]{EG} for details).
\end{definition}
Using Bloch's cycle complex, we can extend this definition to higher Chow groups in a similar manner.

The following definition is a special case of the one in \cite[Section 4.2]{MV}:


\begin{definition}{\cite[Definition 2.1]{krishna2012}}
	A pair $(V,U)$ of smooth schemes over $k$ is said to be a \textit{good pair} for $G$ if $V$ is a  $k$-rational representation of $G$ and $U\subset V$ is a $G$-invariant open subset on which $G$ acts freely and the quotient $U/G$ is a smooth quasi-projective scheme.\\
	A sequence of pairs $\rho=(V_i,U_i)_{i\geq 1}$ is said to be an \textit{admissible gadget} for $G$ if there exists a good pair $(V,U)$ for $G$ such that $V_i=V^{\oplus i}$ and $U_i\subset V_i$ is a $G$-invariant open subscheme such that the following hold:
	\begin{itemize}
		\item $(U_i\oplus V)\cup (V\oplus U_i)\subseteq U_{i+1}$ as $G$-invariant open subsets.
		\item $codim_{U_{i+2}}(U_{i+2}\setminus (U_{i+1}\oplus V))> codim_{U_{i+1}}(U_{i+1}\setminus (U_{i}\oplus V))$.
		\item $codim_{V_{i+1}}(V_{i+1}\setminus U_{i+1})> codim_{V_i}(V_i\setminus U_i)$.
		\item The action of $G$ on $U_i$ is free, and the quotient is quasi-projective.
	\end{itemize}
\end{definition}

An example of such an admissible gadget can be given as follows. Let $V$ be a $k$-rational representation of $GL_n$, and let $U$ be a $GL_r$-invariant open subset on which $GL_r$ acts freely, and the quotient $U/GL_r$ is a quasi-projective scheme. Then $(V,U)$ is a good pair, and we define an admissible gadget $\rho=(V_i,U_i)_{i\geq 1}$ by taking $U_{i+1}:=(U_i\oplus V)\cup (V\oplus U_i)$.\\
For a quasi-projective scheme $X$ with an action $G$, consider the mixed quotients $X^i(\rho):= X\overset{G}{\times} U_i$, and define $X_G(\rho):=\colim X\overset{G}{\times} U_i$. While $
 X\overset{G}{\times} U_i$ is an ordinary colimit, since it is taken over a filtered category, it is also a homotopy colimit (see \cite[Example 12.3.5]{BK}). Further, we denote by $X^{\bullet}_G$ the simplicial scheme:
\begin{center}
\begin{tikzcd}
  \ldots\arrow[r]\arrow[r, shift left]\arrow[r, shift right]\arrow[r, shift right=1.5ex] & G\times G\times X \arrow[r]\arrow[r,shift left]\arrow[r, shift right] & G\times X\arrow[r,shift left]\arrow[r] & X\\	
\end{tikzcd}
\end{center}
The following lemma relates $X_G(\rho)$ with $X^{\bullet}_G$.

\begin{lemma}{\cite[Proposition 3.2]{krishna2012}}
Let $\rho =(V_i,U_i)_{i\geq 1}$ be an admissible gadget for a linear algebraic group $G$ over $k$. For any quasi-projective $G$-scheme $X$, there is a canonical isomorphism $X_G(\rho)\cong X^{\bullet}_G$ in $\mathcal{H}_{\bullet}(k)$.
\end{lemma}


The above lemma gives us a comparison theorem between (higher) Chow groups and motivic cohomology (see also \cite[Theorem 3.5]{Joshua2} for a version in the \'{e}tale topology with $\Q$-coefficients).


%We have put asterisks above instead of actual indices because, as we will see, we get something slightly different from the isomorphisms in Equation \ref{equation-chow-comparison}. Since Definition \ref{totaro-chow} of Chow groups for stacks involves shifting by codimension, the isomorphisms we will also be adjusted by codimension. In fact, we will obtain the following,
%\[CH^j(\sX, 2j- n)\simeq H^{n,i}(\sX,\Z).\]
%Here, $j=i+Nl-r^2$, $l$ is the dimension of the $GL_r$ representation $V$, and $N\gg 0$

\begin{proof}[Proof of Theorem \ref{chow-comparison}]
	Note that $GL_r\times X$ is isomorphic to $X\times_{\sX}X$. Thus, the simplicial scheme $X^{\bullet}_{GL_r}$ is isomorphic to the \v{C}ech hypercover $X_{\bullet}\rightarrow\sX$. By Theorem \ref{motive-construction}, $X_{\bullet}\rightarrow\sX$ is a Nisnevich local equivalence.
	
	Let $\rho=(V_i,U_i)$ be an admissible gadget for $GL_r$, with $\dim (V)=l$, $m:=r^2$. By definition, 
	\[CH^i(\sX, 2i-n)=CH^{i}(X^N(\rho),2i-n) \simeq Hom_{\DM{k,\Z}{}}(M(X^N(\rho)),\Z(i)[n])\] 
	for $N$ sufficiently large. Here, the second equivalence follows from Equation \ref{equation-chow-comparison} and the fact that $H^{n,i}(X,\Z)=Hom_{\DM{k,\Z}{}}(M(X), \Z(i)[n])$. Note that this is well-defined since the maps $X^s(\rho)\rightarrow X^t(\rho)$ are vector bundles and so have the same Chow groups by $\A^1$-invariance. 
	
	Now, as $X_{GL_r}(\rho)=\underset{N}{\colim} X^N(\rho)$ is a filtered colimit, by \cite[Example 12.3.5]{BK} it is also a homotopy colimit. Then, by Remark \ref{remark-motive-commutes-with-homotopy-colimits},
	\begin{align*}
	Hom_{\DM{k,\Z}{}}(M(X_{GL_r}(\rho)), \Z(i)[n])&= 
	Hom_{\DM{k,\Z}{}}(M\big(\underset{N}{\hocolim} X^N(\rho)\big),\Z(i)[n])\\
	&\simeq Hom_{\DM{k,\Z}{}}\big(\underset{N}{\hocolim}M(X^N(\rho)),\Z(i)[n]\big)\\
	&\simeq \underset{N}{\holim}Hom_{\DM{k,\Z}{}}\big(M(X^N(\rho)),\Z(i)[n]\big).
	\end{align*}
	Since the maps $X^n(\rho)\rightarrow X^m(\rho)$ are $\A^1$-invariant, the groups $Hom_{\DM{k,\Z}{}}\big(M(X^N(\rho)),\Z(i)[n]\big)$ stabilise for all $m\geq N$. Thus, 
	\[Hom_{\DM{k,\Z}{}}(M(X_G(\rho)), \Z(i)[n])\simeq Hom_{\DM{k,\Z}{}}\big(M(X^N(\rho)),\Z(i)[n]\big)\]
	whenever $N$ is large enough. Further, we also have the relations,
	\[H^{n,i}(\sX, \Z)=Hom_{\DM{k,\Z}{}}(M(\sX),\Z(i)[n])\simeq Hom_{\DM{k,\Z}{}}(M(X_{\bullet}), \Z(i)[n])\]
	in $\DM{k,\Z}{}$. By the previous lemma, $X_G(\rho)$ and $X_{\bullet}$ are isomorphic in $\mathcal{H}_{\bullet}(k)$. Putting these together, we get required isomorphism.
\end{proof}

\begin{remark}
	Equivariant algebraic cobordisms (see \cite{heller13}, \cite{krishna2012}) are defined by a Borel type construction analogous to definition of equivariant (higher) Chow groups.
	By the above considerations, one can think of equivariant algebraic cobordism as an algebraic cobordism of the associated quotient stack.
\end{remark}




\section{Various Triangles}
\label{section-various-triangles}

We now establish Nisnevich descent and blow-up sequence for the Nisnevich motive and also prove the projective bundle formula. As a consequence of the projective bundle formula, we get a Gysin triangle for cd-quotient stacks. These result are already known for \'{e}tale motives (see \cite{MotivesDM}). All the arguments in this section are directly adapted from their \'{e}tale counterparts in \cite{MotivesDM} $-$ except for the projective bundle formula. The argument for the projective bundle formula in the \'{e}tale case relies on the identification of the Picard group with $H_{\acute{e}t}^2(\sX,\Z(1))$. This identification fails for stacks if \'{e}tale topology is replaced by Nisnevich topology. So we adopt a different approach using homotopical descent.

In this section, we work exclusively with cd-quotient stacks.

\begin{remark}\label{BaseChange}
Let $\sZ$ be a cd-quotient stack. If $\sY\rightarrow\sZ$ is a representable morhpism, then $\sY$ is also a cd-quotient stack. To see this, let $Z\rightarrow [Z/GL_n]\rightarrow\sZ$ be a Nisnevich covering of $\sZ$ by a quotient stack. Now, observe that we have a cartesian diagram by base change,
\begin{center}
\begin{tikzcd}
	\sY\times_{\sZ}Z\arrow[r]\arrow[d] & Z\arrow[d]\\
	\sY\times_{\sZ}[Z/GL_n]\arrow[r]\arrow[d] & \left[ Z/GL_n\right]\arrow[d]\\
	\sY\arrow[r] & \sZ
\end{tikzcd}
\end{center}
where $\sY\times_{\sZ}Z\rightarrow \sY\times_{\sZ} [Z/GL_n]$ is a $GL_n$-torsor and $\sY\times_{\sZ} [Z/GL_n]\rightarrow\sY$ is a Nisnevich cover. Denote the algebraic space $\sY\times_{\sZ}Z$ by $Y$. Thus, $\sY$ is a cd-quotient stack. Hence, by Theorem \ref{motive-construction}, $Y_{\bullet}\rightarrow \sY$ is a Nisnevich local weak equivalence.\\
This tells us that \textit{$GL_n$-presentations} respect base change.
\end{remark}

\begin{proposition}\label{nisnevich-descent} For a distinguished Nisnevich square,
\begin{center}
\begin{tikzcd}
	\sW\arrow[r]\arrow[d] & \sY\arrow[d,"p"]\\
	\sX\arrow[r,hook,"j"]	&\sZ
\end{tikzcd}
\end{center}
where $j$ is an open immersion and $p$ is \'{e}tale representable, the induced diagram on motives is homotopy cartesian.
\end{proposition}

\begin{proof}
For a distinguished Nisnevich square of stacks
\begin{center}
\begin{tikzcd}
	\sW\arrow[r]\arrow[d] & \sY\arrow[d,"p"]\\
	\sX\arrow[r,hook,"j"]	&\sZ
\end{tikzcd}
\end{center}
we need to show that the induced diagram 
\begin{center}
	\begin{tikzcd}
		M(\sZ)\arrow[r,,"p^*"]\arrow[d,"j^*"] & M(\sY)\arrow[d]\\
		M(\sX)\arrow[r,]	&M(\sW)
	\end{tikzcd}
\end{center}
of motives is homotopy cartesian.
If $Z\rightarrow [Z/GL_n]\rightarrow\sZ$ is a Nisnevich covering of $\sZ$ by a $GL_n$-torsor, then by Remark \ref{BaseChange}, we can base change this covering to $\sY,\sX$ and $\sW$. This gives us a cartesian diagram of hypercovers:
\begin{center}
\begin{tikzcd}
	W_{\bullet}\arrow[r]\arrow[d] & Y_{\bullet}\arrow[d,"p_{\bullet}"]\\
	X_{\bullet}\arrow[r,hook,"j_{\bullet}"]	&Z_{\bullet}
\end{tikzcd}
\end{center}
For each $i$, the above the diagram is a distinguished Nisnevich square of schemes. Thus, for each $i$, the following diagram of motives is homotopy (co)cartesian,
\begin{center}
\begin{tikzcd}
	M(Z_i)\arrow[r]\arrow[d] & M(Y_i)\arrow[d]\\
	M(X_i)\arrow[r]	&M(W_i)
\end{tikzcd}
\end{center}
By Remark \ref{remark-motive-commutes-with-homotopy-colimits}, $M(Z_{\bullet})\simeq \hocolim M(Z_i)$. As homotopy colimits commute with homotopy colimits, the following diagram is again homotopy (co)cartesian,
\begin{center}
\begin{tikzcd}
	M(Z_{\bullet})\arrow[r]\arrow[d] & M(Y_{\bullet})\arrow[d]\\
	M(X_{\bullet})\arrow[r]	&M(W_{\bullet})
\end{tikzcd}
\end{center}
By Theorem \ref{motive-construction}, we get the required result.
\end{proof}



\begin{theorem}[Projective Bundle Formula]
\label{projective-bundle-formula}
Let $\sE$ be a vector bundle of rank $n+1$ on a stack $\sX$. There exists a canonical isomorphism in $\DM{k,\Z}{}$:
\[M(\Proj(\sE))\rightarrow\bigoplus_{i=0}^n M(\sX)(i)[2i]\]
\end{theorem}
\begin{proof}
	As projective bundle formula is known for smooth schemes by \cite[Theorem 15.12]{MWV}, we will deduce the result for stacks by a homotopical descent argument. To make such a homotopical descent argument, we need to ensure that homotopy colimits commute with $M$ and derived tensor. But both $M$ and derived tensor are derived functors of left Quillen functors, so Lemma \ref{homotopy-colimits-left-adjoints} ensures that this is true.
	 
	Let $p:\Proj(\sE)\rightarrow \sX$ be the projective bundle, and $\sO(1)$ the canonical line bundle on it. This construction behaves well with respect to base change. If $U_{\bullet}\rightarrow \sX$ is a $GL_n$-presentation, then by base change we get projective bundles $p_i:V_i \rightarrow U_i$ for every $i$, and line bundles $\sO(1)_{V_i}$ on $V_i$ by pullback. Moreover, by Remark \ref{BaseChange}, the \v{C}ech hypercover $V_{\bullet}\rightarrow \Proj(\sE)$ is a $GL_n$-presentation of $\Proj(\sE)$. As each $p_i: V_i\rightarrow U_i$ is a projective bundle, by the projective bundle formula (see \cite[Theorem 15.12]{MWV}), we have:
	
	\begin{equation}
	M(V_i) \simeq \oplus_{j=0}^n M(U_i) \otimes \Z(j)[2j],
	\end{equation}
	in $\DM{k,\Z}{}$.
	
	Since $M$ and derived tensor are left Quillen, using Lemma \ref{homotopy-colimits-left-adjoints}, we have
	\begin{align*}
	M(\Proj(\sE)) \simeq  M( \hocolim V_i )	&\simeq \hocolim (M( V_i ))\\
	&\simeq \hocolim (\oplus_{j=0}^n M( U_i ) \otimes \Z(j)[2j])\\
	&\simeq \oplus_{j=0}^n (M( \hocolim U_i ) \otimes \Z(j)[2j])\\
	&\simeq \oplus_{j=0}^n M(\sX) \otimes \Z(j)[2j],
	\end{align*}
	as required.
	
	
\end{proof}

\begin{remark}
In fact, the above theorem works for any simplicial presheaf $\sX$. This was pointed out to us by a referee. While the notion of a vector bundle is not strightforward to define for simplicial presheaves, one can proceed as follows: for any simplicial presheaf, a principal $GL_n$-bundle is given by a morphism $\sX\rightarrow BGL_n$. In fact, the homotopy fibre product $\sX\times_{BGL_n}^h \Spec k:=\sY$ defines a the total space of the principal $GL_n$-bundle over $\sX$. Let $V$ be an $n$-dimensional vector space over $k$. Now, take the associated vector bundle $\sY\times^{GL_n} V:= \sE$ and projectivise. This gives a projective bundle $\Proj(\sE)\rightarrow \sX$. This construction is compatible with base change.\\
Now by [Dugger, section 2], any simplicial presheaf is quasi-isomorphic to a simplicial presheaf which is levelwise a disjoint union of representables. That is, we have a quasi-isomorphism $X_{\bullet}\rightarrow \sX$, such that each $X_i$ is a disjoint union of schemes. And, the proof above works verbatim.
\end{remark}


\begin{proposition}
	Let $\sZ\subset \sX$ be a smooth closed substack of $\sX$. Let $Bl_{\sZ}(\sX)$ denote the blow-up of $\sX$ in the centre $\sZ$, and $\sE$ be the exceptional divisor. Then we have a canonical distinguised triangle:
	\[M(\sE)\rightarrow M(\sZ)\oplus M(Bl_{\sZ}(\sX))\rightarrow M(\sX)\rightarrow M(\sE)[1]\]
\end{proposition}
\begin{proof}
	Let $X\rightarrow [X/GL_n]\rightarrow\sX$ be a Nisnevich covering of $\sX$ by a $GL_n$-torsor. Since the morphism $Bl_{\sZ}(\sX)\rightarrow\sX$ is projective, it is representable. Then, we can base change $X$ to $Bl_{\sZ}(\sX),\sZ$ and $\sE$. The rest of the proof is the same as Proposition \ref{nisnevich-descent}.
\end{proof}

\begin{theorem}
	Let $\sX$ be a smooth stack and $\sZ \subset \sX$ be a smooth closed substack of pure codimension $c$. Then,
	\[M(Bl_{\sZ}(\sX))\simeq M(\sX)\oplus_{i=0}^{c-1} M(\sZ)(i)[2i]\]
\end{theorem}
\begin{proof}
	Using the previous result, we have a canonical distinguished triangle:
	\[M(\sE)\rightarrow M(\sZ)\oplus M(Bl_{\sZ}(\sX))\rightarrow M(\sX)\rightarrow M(\sE)[1],\]
	where $p:Bl_{\sZ}(\sX)\rightarrow\sX$ is the blow-up. The exceptional divisor is the projectivisation of the normal bundle $N_{\sZ}(\sX)$ of $Z$ in $\sX$, i.e, $\sE\simeq \Proj(N_{\sZ}(\sX))$. If $M(\sX)\rightarrow M(\sE)[1]$ is zero, then the projective bundle formula for $\Proj(N_{\sZ}(\sX))$ gives us the result. 
	
	To prove that $M(\sX)\rightarrow M(\sE)[1]$ is zero, we argue exactly as in \cite[Theorem 3.7]{MotivesDM} (see also \cite[Chapter 5, Proposition 3.5.3]{FSV}). Take $\sX\times \A^1$ and consider the blow-up along $\sZ\times \lbrace 0\rbrace$. We have a map $q:Bl_{\sZ\times\lbrace 0\rbrace}(\sX\times \A^1)\rightarrow \sX\times \A^1$. Consider the morphism of exact triangles,
	\begin{center}
	\begin{tikzcd}
		M(\sE)\arrow[r]\arrow[d] & M(q^{-1}(\sZ\times\lbrace 0\rbrace))\arrow[d]\\
		M(\sZ)\oplus M(Bl_{\sZ}(\sX))\arrow[r]\arrow[d] & M(\sZ\times \lbrace 0\rbrace)\oplus M(Bl_{\sZ\times \lbrace 0\rbrace}(\sX\times \A^1))\arrow[d,"f"]\\
		M(\sX)\arrow[r,"s_0"]\arrow[d,"g"] & M(\sX\times \A^1)\arrow[d,"h"]\\
		M(\sE)[1]\arrow[r,"a"] & M(q^{-1}(\sZ\times\lbrace 0\rbrace))[1]\\
	\end{tikzcd}
	\end{center}
	By the projective bundle formula, the morphism $a$ is split injective, and $s_0$ is an isomorphism. Hence, to show that $g$ is zero, it suffices to show that $h$ is zero. To see this, note that the composition
	\[M(\sX\times \lbrace 1\rbrace)\rightarrow M(Bl_{\sZ\times \lbrace 0\rbrace}(\sX\times \A^1))\rightarrow M(\sX\times \A^1)\]
	is an isomorphism. This implies that $f$ admits a section so $h$ must be zero.
\end{proof}

\begin{definition}
For a map $M(X)\rightarrow M(Y)$ of motives of stacks (or simplicial schemes), we denote the cone by
\[M\Big(\frac{X}{Y}\Big):= cone(M(X)\rightarrow M(Y))\;\; \text{in} \;\DM{k,\Z}{}.\]
\end{definition}

\begin{lemma}
	Let $f:\sX'\rightarrow\sX$ be an \'{e}tale representable morphism of algebraic stacks, and let $\sZ\subset \sX$ be a closed substack such that $f$ induces an isomorphism $f^{-1}(\sZ)\simeq \sZ$. Then the canonical morphism
	\[M\Big(\frac{\sX'}{\sX'\setminus\sZ}\Big)\rightarrow M\Big(\frac{\sX}{\sX\setminus\sZ}\Big)\]
	is an isomorphism.
\end{lemma}
\begin{proof}
	Let $U_{\bullet}\rightarrow\sX$ be a $GL_n$-presentation. Let $V_{\bullet}:=U_{\bullet}\times_{\sX} \sX'$ be a $GL_n$-presentation of $\sX'$ obtained by base change. We have a map of simplicial sets $f_{\bullet}:V_{\bullet}\rightarrow U_{\bullet}$ induced by the map $f:\sX' \rightarrow \sX$. Note that for every simplicial degree, $f_i$ is \'{e}tale and that $f_i^{-1}(\sZ\times_{\sX} U_i)\simeq \sZ\times_{\sX} U_i$. Let $Z_i$ denote the base change $\sZ\times_{\sX}U_i$. Then, by \cite[Chapter 3, Proposition 5.18]{FSV}, the canonical morphism
	\[M\Big(\frac{V_{\bullet}}{V_{\bullet}\setminus Z_{\bullet}}\Big)\rightarrow M\Big(\frac{U_{\bullet}}{U_{\bullet}\setminus Z_{\bullet}}\Big).\]
	is an isomorphism. Now, Theorem \ref{motive-construction} gives us the required result.
\end{proof}


\begin{corollary}
	Let $p:V\rightarrow\sX$ be a vector bundle of rank $n$ over an algebraic stack. Denote by $s:\sX\rightarrow V$ the zero section. Then
	\[M\Big(\frac{V}{V\setminus s}\Big)\simeq M(\sX)(d)[2d].\]
\end{corollary}
\begin{proof}
	From the previous lemma we have an isomorphism
	\[M\Big(\frac{V}{V\setminus s}\Big)\simeq M\Big(\frac{\Proj(V\oplus \sO)}{\Proj(V\oplus \sO\setminus s)}\Big),\]
	and we have 
	\[M\Big(\frac{\Proj(V\oplus \sO)}{\Proj(V\oplus \sO\setminus s)}\Big)\simeq M(\sX)(d)[2d]\]
	from the projective bundle formula (see \cite[Lemma 3.9]{MotivesDM} for details).
\end{proof}

\begin{theorem}[Gysin Triangle]
	\label{gysin-triangle}
	Let $\sZ \subset \sX$ be a smooth closed substack of codimension c. Then there exists a Gysin triangle:
	\[M(\sX \setminus \sZ)\rightarrow M(\sX)\rightarrow M(\sZ)(c)[2c] \rightarrow M(\sX \setminus \sZ)[1].\]
\end{theorem}
\begin{proof}
	Note that we have an exact triangle
	\[M(\sX \setminus \sZ)\rightarrow M(\sX)\rightarrow M\Big(\frac{\sX}{\sX\setminus\sZ}\Big) \rightarrow M(\sX \setminus \sZ)[1].\]
	So it suffices to show that $M\big(\frac{\sX}{\sX\setminus\sZ}\big)\simeq M(Z)(c)[2c]$ in $\DM{k,\Z}{}$. The argument is exactly as in \cite[Theorem 3.10]{MotivesDM}.
\end{proof}

\section{Hypercohomology}

\begin{definition}[Smooth-Nisnevich site]
	Let $\sX$ be an algebraic stack. The smooth-Nisnevich site of $\sX$, denoted by $\sX_{sN}$ is the category whose objects consist of pairs $(U,p)$ where $p$ is a smooth morphism $p:U\rightarrow \sX$  from an algebraic space $U$. Coverings in $\sX_{sN}$ are given by Nisnevich covering of algebraic spaces.
\end{definition}


\begin{theorem}
	The motivic cohomology of $\sX$ agrees with the hypercohomology of the constant sheaf $\Z$ on $\sX_{sN}$. That is,
	\[Ext^i_{DA(\sX)}(\Z,\Z(j))\simeq Hom_{\DM{k,\Z}{}}(M(\sX),\Z(j)[i]),\]
	where $\Z$ denotes the constant sheaf $\Z$ on the $\sX_{sN}$.
\end{theorem}

\begin{proof}
	We will us the fact that $\sX$ has a $GL_n$-presentation $p: U_{\bullet}\rightarrow\sX$. The morphism $p: U_{\bullet}\rightarrow \sX$ is a local weak equivalence over $\sX$, and $\Z(U_{\bullet})\simeq \Z$ in $DA(\sX)$. Thus, we have
	\[Ext^i(\Z,\Z(j))= Ext^i_{DA(\sX)}(\Z(U_{\bullet}),\Z(j)).\]
	Let $g:U_{\bullet}\rightarrow \Spec k$ and $h: \sX\rightarrow \Spec k$ be the structure maps. Since, $\Z(j)$ in $DA(\sX)$ is the pullback of motivic sheaf $\Z(j)$ over $k$, we have,
	
	\begin{align*}
	Ext^i_{DA(\sX)}	(\Z(U_{\bullet}),h^*\Z(j)) &= Ext^i_{DA(U_{\bullet})}(p^*\Z(U_{\bullet}),p^*h^*\Z(j))\\
	&= Ext^i_{DA(k)}(g_\#(\Z|_{U_{\bullet}}),\Z(j))\\
	&= Ext^i_{DA(k)}(\Z(U_{\bullet}),\Z(j)).
	\end{align*}
	Since $\Z(j)$ are $\A^1$-local complexes, we have that
	\[Ext^i(\Z(U_{\bullet}),\Z(j)) = Hom_{\DM{k,\Z}{}}(M(U_{\bullet}),\Z(j)[i]).\]
	
	Now, the result follows since $U_{\bullet}\rightarrow \sX$ is a local weak equivalence in $\sH(k)$.
\end{proof}


\section{Application to Exhaustive stacks}
\label{section-exhaustive}
In \cite{hoskins2019}, a motive $M_{exh}$ is defined for a class of stacks which they call as \textit{exhaustive stacks}. They do this by using an idea similar to Totaro's ``finite-dimensional approximation" technique in \cite{TotaroChow}. Examples of exhaustive stacks are quotient stacks and the moduli stack of vector bundles on a curve of fixed rank and degree. In fact, exhaustive stacks turn out to be special cases of cd-quotient stacks (see Lemma \ref{proposition-filtration-by-quotients}). We will compare the motive $M_{exh}$ with the motive of Definition \ref{definition-motive-stack}.

In this section, we adopt the conventions used in \cite{hoskins2019} for algebraic stacks. In particular this means that we work with stacks $\sX$ which admit a smooth atlas $p:U\rightarrow \sX$ by a locally finite type $k$-scheme $U$ such that $p$ is schematic (representable by \textit{schemes})\footnote{This is only to maintain consistency with the conventions in \cite{hoskins2019}. It does not particularly affect the arguments that we present, which work for any stack locally of finite type over $k$.}.


\begin{definition}[\!\! {\cite[Definition 2.15]{hoskins2019}}]
	Let $\sX$ be an algebraic stack locally of finite type over a field $k$. Let $\sX_0 \subset \sX_1\subset \ldots$ be an increasing filtration of $\sX$ such $\sX_i\subset \sX$ are quasi-compact open substacks and their union covers $\sX$, i.e, $\sX=\cup_i \sX_i$. Then an \textit{exhaustive sequence of vector bundles} with respect to this filtation is a sequence of pairs $\{(V_i,W_i)\}_{i\geq 0}$ where $V_i$ is a vector bundle on $\sX_i$ and $W_i\subset V_i$ is a closed substack such that
	\begin{enumerate}
		\item the complement $U_i:= V_i\setminus W_i$ is a separated $k$-scheme of finite type,
		\item we have injective maps of vector bundles $f_{i,i+1}:V_i\rightarrow V_{i+1}\times_{\sX_{i+1}}\sX_i$ such that $f_{i,i+1}^{-1}(W_{i+1}\times_{\sX_{i+1}}\sX_i)\subset W_i$ and,
		\item the codimension of $W_i$ in $V_i$ tends to infinity as $i$ increases.
	\end{enumerate}
	A stack admitting an exhaustive sequence with respect to some filtration is said to be \textit{exhaustive}.
\end{definition}


In fact, one can show that every exhaustive stack admits a filtration by global quotient stacks. This implies that it is a cd-quotient stack.

\begin{lemma}
	\label{proposition-filtration-by-quotients}
	Let $\sX$ be an exhaustive stack. Let $\sX=\cup_i\sX_i$  be an increasing filtration with an exhaustive sequence of vector bundles. Then there exists an increasing filtration $\sX=\cup_i\sY_i$ with $\sY_i\subseteq \sX_i$ and each $\sY_i$ is a global quotient stack. In particular, it is a cd-quotient stack.
\end{lemma}

\begin{proof}
	Let $\{(V_{i},W_i)\}_{i\geq 0}$ denote the exhaustive sequence of vector bundles corresponding to the filtration $\{\sX_i\}_{i\geq 0}$. Let $p_i:V_i\rightarrow \sX_i$ denote the structure map of the vector bundle $V_i$. Now by definition the complement $U_i= V_i\setminus W_i$ is a separated finite type $k$-scheme. Since $p_i$ is smooth, the image $p_i(U_i)\subset \sX_i$ is an open substack which is of finite type over $k$. Set $\sY_i:=p_i(U_i)$. Consider the restriction $V_i\times_{\sX_i}\sY_i\rightarrow \sY_i$ of $V_i$ to this substack. This is a vector bundle on $\sY_i$ which contains an open representable substack $U_i\subset V_i\times_{\sX_i}\sY_i$ that surjects onto $\sY_i$. Thus, $\sY_i$ is a global quotient stack, by \cite[Lemma 2.12]{EHKV}.
	
	Thus, we have an increasing filtration $\{\sY_i\}_{i\geq 0}$ such that $\sY_i\subseteq \sX_i$. The only thing left to check is that this filtration covers $\sX$. This follows from the following topological argument.
	
	Take a point $x\in \sX$. Then as the filtration $\{\sX_i\}_{i\geq 0}$ covers $\sX$, there exists an $i$ such that $x\in \sX_i$. We will show that there exists an $N\geq i$ such that $x\in \sY_N$. 
	
	If $x\in \sY_i$, there is nothing to prove. So assume that $x\notin \sY_i$. This means that $p_i^{-1}(x)\subset W_i$ in the vector bundle $V_i$. Let $Z:=\overline{p_i^{-1}\{x\}}$ be the closure of the fibre in $V_i$. Since $Z\subseteq W_i$ we see that 
	\[n:=\codim\, Z\geq \codim W_i.\]
	As $\{(V_i,W_i)\}_{i\geq 0}$ is an exhaustive sequence, there exists an $N$ such that $\codim\, W_N > n$. Further, we have a map $f_{i,N}:V_i\rightarrow V_N$ such that $f_{i,N}^{-1}(W_N)\subset W_i$. If $Z$ was contained in $f_{i,N}^{-1}(W_N)$, we would have
	\[n=\codim\, Z\geq \codim \, f_{i,N}^{-1}(W_N)> n,\]
	a contradiction. Thus, there exists $y\in p_i^{-1}(x)$ such that $f_{i,N}(y)\in U_N$ implying that $x\in \sY_N$.
\end{proof}


\begin{definition}[\!\!{\cite[Definition 2.17]{hoskins2019}}]
Let $\sX$ be an exhaustive stack wiht an exhaustive sequence of vector bundles $\{(V_i,W_i)\}_{i\geq 0}$. The motive $M_{exh}(\sX)$ is defined in as the colimit of the motives of the schemes $U_i$. That is,
\[M_{exh}(\sX)=\colim M(U_i).\]
\end{definition}

Since exhaustive stacks are cd-quotient stacks, we would like to compare the motive $M_{exh}(\sX)$ with the motive $M(\sX)$ in Definition \ref{definition-motive-stack}. The following proposition shows that they are isomorphic in $\DM{k,\Z}{}$.


\begin{proposition}
	\label{proposition-comparision-with-hoskins19}
	Let $\sX$ be a smooth exhaustive stack. Then
	\[M(\sX) \simeq M_{exh}(\sX) \;\; \text{in}\;\; \DM{k,\Z}{}\]
\end{proposition}
\begin{proof}
	Let $X\rightarrow \sX$ be the $0$-skeleton of a $GL_n$-presentation. By \cite[Theorem II.6.4]{knutson}, there exists a Nisnevich covering $Y\rightarrow X$ with $Y$ a scheme. This gives us a presentation $Y\rightarrow \sX$. Let $Y_{\bullet}\rightarrow \sX$ be the associated \v{C}ech hypercover. By similar argument as in Theorem \ref{motive-construction}, we get that $Y_{\bullet} \simeq \sX$ in $\mathcal{H}_{\bullet}(k)$, and so we have $M(Y_{\bullet}) \simeq M(\sX)$ in $\DM{k,\Z}{}$. Hence, it suffices to show that $M(Y_{\bullet}) \simeq M_{exh}(\sX)$ in $\DM{k,\Z}{}$.	The proof is exactly the same as \cite[Proposition A.7]{hoskins2019} using the atlas $Y_{\bullet}\rightarrow\sX$.
\end{proof}


\begin{example}
	Let $C$ be a smooth projective geometrically connected curve of genus $g$ over a field $k$.	In \cite[Section 3]{hoskins2019}, it is proved that the moduli stack $Bun_{n,d}$ of vector bundles on a $C$ of fixed rank $n$ and degree $d$ is exhaustive. To show this, they observe that it admits a filtration by the maximal slope of all vector bundles. Thus, $Bun_{n,d}$ is a filtered colimit of open substacks $Bun_{n,d}^{\geq \mu_l}$ where $\{\mu_l\}$ is an increasing sequence of rational number representing the maximal slope. Also, each of these open substacks is a global quotient stack and can be written as $Bun_{n,d}^{\geq \mu_l}:=[Q^{\geq \mu_l}/GL_N]$ where $Q^{\geq \mu_l}$ is an open subscheme of a Quot scheme (see \cite[Th\'{e}or\`{e}me 4.6.2.1]{LMB} for futher details). Then, by \cite[Lemma 2.26]{hoskins2019}, we have
	\[M_{exh}(Bun_{n,d})=\hocolim_l M_{exh}(Bun_{n,d}^{\geq \mu_l})\]
	and from Proposition \ref{proposition-comparision-with-hoskins19} we get that
	\[M(Bun_{n,d})=\hocolim_l M(Bun_{n,d}^{\geq \mu_l})=\hocolim_l M(Q^{\geq \mu_l}_{\bullet}).\]
	Thus, for $Bun_{n,d}$, the motive $M_{exh}$ of \cite{hoskins2019} can be computed as a homotopy colimit of the motive of Definition \ref{definition-motive-stack}.
	Since these homotopy colimits are being taken over filtered categories, they can actually be computed by their ordinary colimits (see \cite[Example 12.3.5]{BK}).
\end{example}

\begin{remark}
	Proposition \ref{proposition-comparision-with-hoskins19} shows that for exhaustive stacks the motive defined using the nerve construction in Definition \ref{definition-motive-stack} agrees with $M_{exh}$ defined in \cite{hoskins2019} in the Nisnevich topology. This potentially simplifies many of the functoriality arguments in \cite[\S 2]{hoskins2019}. For example, it is now immediate that $M_{exh}$ is independent of the choices involved in its construction (see also \cite[Lemma 2.20]{hoskins2019}). In \cite[Appendix A]{hoskins2019}, such a comparison is proved for \'{e}tale motives.
\end{remark}






\appendix
\section{Homotopy (co)limits and derived functors}



Let $I$ be an index category, and let $M^I$ denote the category of $I$-diagrams in $M$. $M^I$ is just the cateogry of functors $Fun(I,M)$. We have a constant functor $c:M\rightarrow M^I$, taking every object $A$ of $M$ to the constant $I$-diagram whose every object is $A$ and all maps are $id_A$. The left adjoint of $c$ (if it exists) is the colimit functor $\colim :M^I\rightarrow M$.\\
Let $F:M\rightleftarrows N: G$ be an adjunction. Then we also have an induced adjunction between the category of $I$-diagrams,
\[F^I:M^I\rightleftarrows N^I:G.\]

The following lemma shows that just as left adjoints commute with colimits, derived functors of left adjoints commute with \textit{homotopy colimits}.

\begin{lemma}\label{homotopy-colimits-left-adjoints}
	Let $F:M\rightleftarrows N: G$ be a Quillen adjunction of model categories. Let $I$ be an index category such that the projective model structure is defined on $M^I$ and $N^I$. Then, for an $I$-diagram $E$ in $M$, we have
	\[LF(\hocolim E)\simeq \hocolim LF^I(E),\]
	where $LF$ and $LF^I$ are the derived functors of $F$ and $F^I$, repectively.
\end{lemma}
\begin{proof}
	Note that $LF$ is defined as the composite
	\[Ho(M)\overset{Q}{\rightarrow} Ho(M_c)\overset{F}{\rightarrow} Ho(N)\]
	where $Q$ is the cofibrant replacement functor.
	
	Let $M^I$ denote the category of $I$-diagrams in $M$. Note that the fibrations and weak equivalences in the projective model structure on $M^I$ are defined object-wise. The homotopy colimit functor is the derived functor of the colimit functor $\colim : M^I\rightarrow M$ which is the left adjoint of the constant functor $M\rightarrow M^I$. More precisely,
	\[\hocolim: Ho(M^I)\overset{Q}{\rightarrow}Ho(M^I)\overset{\colim}{\rightarrow}Ho(M).\]
	Note that here $Q$ is the cofibrant replacement functor in the \textit{projective model structure} of $M^I$.
	This says that for any $I$-diagram $E$ the homotopy colimit can be computed by taking the ordinary colimit of its cofibrant replacement $QE$, i.e,
	\[\hocolim E\simeq \colim QE\]
	is a weak equivalence in the homotopy category.
	Since, $M^I$ has the projective model structure, {\rm colim} is left Quillen. Thus, $\colim QE$ is, in fact, a cofibrant object in $M$ and
	\begin{align*}
	LF(\hocolim E) &\simeq LF(\colim QE) \\
	&\simeq F(\colim QE).
	\end{align*}
	
	Now, observe that the adjoint pair $(F,G)$ induces an adjunction on the diagram categories associated to $I$,
	\[F^I:M^I\rightleftarrows N^I:G^I.\]
	As fibrations are defined object-wise and $G$ is right Quillen, $G^I$ preserves fibrations. Hence, $F^I$ is left Quillen, and preserves cofibrations. This means that the image $F^I(QE)$ is cofibrant in $N^I$.
	Then, we have
	\begin{align*}
	\hocolim LF^I(E) &\simeq \hocolim F^I(QE)\\
	&\simeq \colim F^I(QE).
	\end{align*}

As $F$ is a left adjoint, it commutes with ordinary colimits. That is, we have a commutative diagram,
\begin{center}
	\begin{tikzcd}
	M^I\arrow[r,"F^I"]\arrow[d] & N^I\arrow[d]\\
	M\arrow[r,"F"] &N
	\end{tikzcd}
\end{center}
where the vertical arrows are the colimit functor. Thus,
\[F(\colim QE)=\colim F^I(QE),\]
as required.
\end{proof}

\begin{remark}
	Except in some special cases (for example, filtered colimits), homotopy (co)limits and ordinary (co)limits almost never agree. 
	One can greatly simply the discussions (and the confusions arising from it) about ordinary and homotopy (co)limits by using the language of $\infty$-categories. (Co)Limits in $\infty$-categories are defined in such a way that they give the correct notion of homotopy (co)limits. In particular, there are only homotopy colimits in the world of $\infty$-categories. Moreover, when restricted to ordinary categories they recover the usual notion of (co)limits.\\
	It is proved in \cite[Proposition 5.2.3.5]{HTT} that colimits (in $\infty$-categories) commute with left adjoint functors between $\infty$-categories (and limits commute with right adjoints). By \cite{DKhammock,DKsimploc}, all model categories admit a simplicial localisation. This simplicial localisation can be viewed as an $\infty$-category (see \cite{HTT}). So the above lemma may be seen as an alternative version (or a special case, depending on the reader's perspective) of \cite[Proposition 5.2.3.5]{HTT}.
\end{remark}


\end{comment}

%\vspace{1 cm}
%\Green{Things to add:
%\begin{enumerate}
%	\item Chow groups on tensoring with $\Q$.
%	\item comparison with Dhillon-Behrend (?).
%\end{enumerate}
%}


%%%%%%%%%%%%%%%%%%%%%%%

\bibliography{mybib-JACK-EDIT.bib}
\bibliographystyle{amsalpha}



\end{document}
